%% toc: P3 35 35.1,35.2,35.3,35.4
\section{Proofs for comparison games and infinitary languages}
\label{sec:p03-w03-p03-038-annex}

Throughout, \(\tau\) is relational, structures have nonempty domains, and
for a tuple \(\bar a=(a_1,\ldots,a_n)\) and a map
\(\pi:\{1,\ldots,m\}\to\{1,\ldots,n\}\) we write
\(\bar a\circ\pi\defeq(a_{\pi(1)},\ldots,a_{\pi(m)})\). Two facts about
satisfaction are used without further comment, each proved by induction on
formulas from the satisfaction clauses: satisfaction of a formula depends
only on the values of its free variables, and satisfaction of a formula in
a subsignature \(\tau_0\) depends only on the interpretations of the symbols
of \(\tau_0\).

\begin{annexitem}{annex:p03-w03-p03-038-finite-records}{thm:p03-w03-p03-038-finite-records}{Finitely many records}
\begin{proof}
The list \(\mathrm{At}_n\) has \(n^2\) equality formulas and
\(n^{r}\) formulas \(R(x_{i_1},\ldots,x_{i_r})\) for each \(r\)-ary
\(R\in\tau\). As \(\tau\) is finite,
\(N_n=n^2+\sum_{R\in\tau}n^{\operatorname{ar}(R)}\) is finite. We prove by
induction on \(j\), simultaneously for all \(n\), that
\(\mathrm{Rec}_{n,j}\) is finite with the stated size. For \(j=0\),
\(\mathrm{Rec}_{n,0}=\mathcal P(\mathrm{At}_n)\) has \(2^{N_n}\) members.
If \(\mathrm{Rec}_{n+1,j}\) is finite, then
\(\mathrm{Rec}_{n,j+1}=\mathcal P(\mathrm{At}_n)\times
\mathcal P(\mathrm{Rec}_{n+1,j})\) is a product of two finite sets with
\(2^{N_n}\cdot2^{|\mathrm{Rec}_{n+1,j}|}\) members.
\end{proof}
\end{annexitem}

\begin{annexitem}{annex:p03-w03-p03-038-characteristic}{thm:p03-w03-p03-038-characteristic}{Characteristic formulas}
\begin{proof}
All claims are proved by induction on \(j\), simultaneously for all \(n\)
and all \(t\in\mathrm{Rec}_{n,j}\).

\emph{Free variables and rank.} For \(j=0\), \(\chi_S\) is a finite
conjunction of members of \(\mathrm{At}_n\) and their negations, so its
free variables are among \(x_1,\ldots,x_n\) and its rank is \(0\); the empty
conjunction \(\neg\bottom\) also has rank \(0\). For \(t=(S,Q)\) with
\(Q\subseteq\mathrm{Rec}_{n+1,j}\), each \(\chi_q\) has free variables among
\(x_1,\ldots,x_{n+1}\) and rank at most \(j\) by the induction hypothesis.
Quantifying \(x_{n+1}\) leaves free variables among \(x_1,\ldots,x_n\) and
raises the rank by one; a finite conjunction or disjunction has the maximum
rank of its members, and \(\bottom\) has rank \(0\). Hence
\(\operatorname{qr}(\chi_t)\le j+1\).

\emph{Definability, case \(j=0\).} \(\mathcal A\models\chi_S[\bar a]\) iff
every \(\alpha\in S\) holds at \(\bar a\) and every
\(\alpha\in\mathrm{At}_n\setminus S\) fails there, iff
\(\{\alpha\in\mathrm{At}_n:\mathcal A\models\alpha[\bar a]\}=S\), iff
\(\operatorname{rec}_0(\mathcal A,\bar a)=S\).

\emph{Definability, case \(j+1\).} Let \(t=(S,Q)\) and put
\(Q'\defeq\{\operatorname{rec}_j(\mathcal A,\bar ac):c\in\dom\mathcal A\}\),
so that \(\operatorname{rec}_{j+1}(\mathcal A,\bar a)=
(\operatorname{rec}_0(\mathcal A,\bar a),Q')\). For \(c\in\dom\mathcal A\)
and \(q\in\mathrm{Rec}_{n+1,j}\), the formula
\(\chi_q(x_1,\ldots,x_{n+1})\) holds under the assignment
\(x_i\mapsto a_i\), \(x_{n+1}\mapsto c\) iff \(\mathcal A\models\chi_q[\bar ac]\),
iff \(\operatorname{rec}_j(\mathcal A,\bar ac)=q\), by the induction
hypothesis.

Suppose \(\operatorname{rec}_{j+1}(\mathcal A,\bar a)=t\), so
\(\operatorname{rec}_0(\mathcal A,\bar a)=S\) and \(Q'=Q\). The conjunct
\(\chi_S\) holds by the case \(j=0\). For \(q\in Q=Q'\) choose \(c\) with
\(\operatorname{rec}_j(\mathcal A,\bar ac)=q\); then \(\chi_q\) holds at
\(\bar ac\), so \(\exists x_{n+1}\chi_q\) holds at \(\bar a\). For every
\(c\), \(q_c\defeq\operatorname{rec}_j(\mathcal A,\bar ac)\in Q\) and
\(\chi_{q_c}\) holds at \(\bar ac\), so the disjunction over \(Q\) holds at
\(\bar ac\); hence the universal conjunct holds. Thus
\(\mathcal A\models\chi_t[\bar a]\).

Conversely suppose \(\mathcal A\models\chi_t[\bar a]\). The conjunct
\(\chi_S\) gives \(\operatorname{rec}_0(\mathcal A,\bar a)=S\). If \(q\in Q\),
the conjunct \(\exists x_{n+1}\chi_q\) supplies \(c\) with \(\chi_q\) true
at \(\bar ac\), whence \(q=\operatorname{rec}_j(\mathcal A,\bar ac)\in Q'\).
If \(c\in\dom\mathcal A\), the universal conjunct supplies \(q\in Q\) with
\(\chi_q\) true at \(\bar ac\), whence
\(\operatorname{rec}_j(\mathcal A,\bar ac)=q\in Q\); so \(Q'\subseteq Q\).
Therefore \(Q'=Q\) and \(\operatorname{rec}_{j+1}(\mathcal A,\bar a)=(S,Q)=t\).
Records with \(Q=\varnothing\) are covered: then \(\chi_t\) contains
\(\forall x_{n+1}\bottom\), which no structure with a nonempty domain
satisfies, and no record of a tuple in such a structure has an empty second
component.
\end{proof}
\end{annexitem}

\begin{annexitem}{annex:p03-w03-p03-038-pointed-ef}{thm:p03-w03-p03-038-pointed-ef}{The finite game theorem with parameters}
\begin{proof}
\emph{Projection lemma.} If
\(\operatorname{rec}_j(\mathcal A,\bar a)=\operatorname{rec}_j(\mathcal B,\bar b)\)
for tuples of length \(n\), then
\(\operatorname{rec}_j(\mathcal A,\bar a\circ\pi)=
\operatorname{rec}_j(\mathcal B,\bar b\circ\pi)\) for every
\(\pi:\{1,\ldots,m\}\to\{1,\ldots,n\}\). For \(\alpha\in\mathrm{At}_m\) let
\(\alpha^\pi\in\mathrm{At}_n\) replace each \(x_i\) by \(x_{\pi(i)}\). By the
atomic satisfaction clause,
\(\mathcal A\models\alpha[\bar a\circ\pi]\) iff
\(\mathcal A\models\alpha^\pi[\bar a]\). Hence
\(\operatorname{rec}_0(\mathcal A,\bar a\circ\pi)=
\{\alpha:\alpha^\pi\in\operatorname{rec}_0(\mathcal A,\bar a)\}\) is computed
from \(\operatorname{rec}_0(\mathcal A,\bar a)\), the first component of every
record of \(\bar a\); this settles \(j=0\) and the first component for all
\(j\). For \(j+1\), let \(\pi'\) extend \(\pi\) by \(m+1\mapsto n+1\), so that
\((\bar a\circ\pi)c=(\bar ac)\circ\pi'\). Given \(c\in\dom\mathcal A\), equality
of the \((j+1)\)-records gives \(d\) with
\(\operatorname{rec}_j(\mathcal A,\bar ac)=\operatorname{rec}_j(\mathcal B,\bar bd)\),
and the induction hypothesis applied to \(\pi'\) gives
\(\operatorname{rec}_j(\mathcal A,(\bar a\circ\pi)c)=
\operatorname{rec}_j(\mathcal B,(\bar b\circ\pi)d)\). So every extension record
of \(\bar a\circ\pi\) is one of \(\bar b\circ\pi\), and the symmetric argument
gives the converse inclusion.

\emph{(1) implies (2).} Induction on \(k\). For \(k=0\) the only play
consists of the initial position, so a winning strategy exists only if
\((\bar a,\bar b)\) is a partial isomorphism, i.e.\
\(\operatorname{rec}_0(\mathcal A,\bar a)=\operatorname{rec}_0(\mathcal B,\bar b)\).
Let \(\sigma\) win \(G_{k+1}(\mathcal A,\bar a;\mathcal B,\bar b)\). Domains are
nonempty, so plays exist; the final position of each extends
\((\bar a,\bar b)\) and is a partial isomorphism, and its restriction to the
first \(n\) coordinates is one too. Hence the \(0\)-records agree. For
\(c\in\dom\mathcal A\), let \(d\) be \(\sigma\)'s answer to the move \(c\).
The continuation of \(\sigma\) after this round wins
\(G_k(\mathcal A,\bar ac;\mathcal B,\bar bd)\), and the induction hypothesis
gives \(\operatorname{rec}_k(\mathcal A,\bar ac)=\operatorname{rec}_k(\mathcal B,\bar bd)\).
Moves \(d\in\dom\mathcal B\) are treated symmetrically. Thus the extension
record sets coincide, and the \((k+1)\)-records are equal.

\emph{(2) implies (1).} Induction on \(k\). For \(k=0\), equal \(0\)-records
mean that the initial position is a partial isomorphism, which is the final
position of the unique play. For \(k+1\), equality of the extension record
sets means that for every \(c\in\dom\mathcal A\) the set of
\(d\in\dom\mathcal B\) with
\(\operatorname{rec}_k(\mathcal A,\bar ac)=\operatorname{rec}_k(\mathcal B,\bar bd)\)
is nonempty, and symmetrically. By the Axiom of Choice fix such an answer
\(d_c\) for each \(c\), an answer \(c_d\) for each \(d\), and, by the
induction hypothesis, a winning strategy for
\(G_k\) from each resulting position. Duplicator answers the first move with
the fixed answer and then follows the fixed strategy; every play ends in a
partial isomorphism.

\emph{(2) implies (3).} We prove: for every formula \(\varphi\), every \(n\),
every list \(\bar v=(v_1,\ldots,v_n)\) of distinct variables containing the
free variables of \(\varphi\), every \(k\ge\operatorname{qr}(\varphi)\), and all
\(\bar a,\bar b\) of length \(n\) with equal \(k\)-records,
\(\mathcal A\models\varphi[\bar v:=\bar a]\) iff
\(\mathcal B\models\varphi[\bar v:=\bar b]\). Here \([\bar v:=\bar a]\) is any
assignment sending \(v_i\) to \(a_i\). The induction is on \(\varphi\).

If \(\varphi\) is \(\bottom\) both sides fail. If \(\varphi\) is
\(v_i=v_j\) or \(R(v_{i_1},\ldots,v_{i_r})\), it holds at \(\bar a\) iff the
member \(x_i=x_j\), respectively \(R(x_{i_1},\ldots,x_{i_r})\), of
\(\mathrm{At}_n\) lies in \(\operatorname{rec}_0(\mathcal A,\bar a)\), the
first component of the \(k\)-record; the same holds for \(\bar b\). Negation
and the binary connectives follow from the induction hypothesis with the same
\(k\), since their arguments have rank at most \(k\).

Let \(\varphi=\exists y\,\psi\) and \(k=k'+1\) with \(k'\ge\operatorname{qr}(\psi)\).
Let \(\pi\) enumerate in increasing order the indices \(i\) with \(v_i\ne y\),
say \(m\) of them, and let \(\bar v'\defeq(v_{\pi(1)},\ldots,v_{\pi(m)},y)\).
These variables are distinct and include every free variable of \(\psi\). As
\(y\) is not free in \(\varphi\), \(\mathcal A\models\varphi[\bar v:=\bar a]\)
iff there is \(c\) with \(\mathcal A\models\psi[\bar v':=(\bar a\circ\pi)c]\).
By the projection lemma the \(k\)-records of \(\bar a\circ\pi\) and
\(\bar b\circ\pi\) agree. Given such \(c\), choose \(d\) with
\(\operatorname{rec}_{k'}(\mathcal A,(\bar a\circ\pi)c)=
\operatorname{rec}_{k'}(\mathcal B,(\bar b\circ\pi)d)\); the induction
hypothesis for \(\psi\), \(\bar v'\) and \(k'\) gives
\(\mathcal B\models\psi[\bar v':=(\bar b\circ\pi)d]\), so
\(\mathcal B\models\varphi[\bar v:=\bar b]\). The converse is symmetric. If
\(\varphi=\forall y\,\psi\), define \(\pi,\bar v',k'\) in the same way. Suppose
\(\mathcal A\models\varphi[\bar v:=\bar a]\) and let \(d\in\dom\mathcal B\).
The back half of record equality gives \(c\) with equal \(k'\)-records of
\((\bar a\circ\pi)c\) and \((\bar b\circ\pi)d\); as \(\psi\) holds at
\((\bar a\circ\pi)c\), the induction hypothesis makes it hold at
\((\bar b\circ\pi)d\). So \(\mathcal B\models\varphi[\bar v:=\bar b]\), and the
converse is symmetric. Taking \(\bar v=(x_1,\ldots,x_n)\) gives (3).

\emph{(3) implies (2).} Let \(t=\operatorname{rec}_k(\mathcal A,\bar a)\). By
\cref{thm:p03-w03-p03-038-characteristic}, \(\chi_t\) has rank at most \(k\),
free variables among \(x_1,\ldots,x_n\), and holds at \(\bar a\). By (3) it
holds at \(\bar b\), and the same lemma gives
\(\operatorname{rec}_k(\mathcal B,\bar b)=t\).

The finiteness of \(\tau\) entered only through the existence of \(\chi_t\);
the proofs of (1) implies (2) and (2) implies (3) apply verbatim to any
relational signature. For \(n=0\) the theorem is \cref{thm:p03-ef}.
\end{proof}
\end{annexitem}

\begin{annexitem}{annex:p03-w03-p03-038-finitely-many}{thm:p03-w03-p03-038-finitely-many}{Finitely many formulas of each rank}
\begin{proof}
Let \(\varphi\) have rank at most \(k\) and free variables among
\(x_1,\ldots,x_n\), and let \(X\) be the set of \(t\in\mathrm{Rec}_{n,k}\)
such that some structure \(\mathcal A'\) and tuple \(\bar a'\) have
\(\operatorname{rec}_k(\mathcal A',\bar a')=t\) and
\(\mathcal A'\models\varphi[\bar a']\). If \(\mathcal A\models\varphi[\bar a]\),
then \(t=\operatorname{rec}_k(\mathcal A,\bar a)\in X\) and \(\chi_t\) holds at
\(\bar a\). If \(\chi_t\) holds at \(\bar a\) for some \(t\in X\), then
\(\operatorname{rec}_k(\mathcal A,\bar a)=t=\operatorname{rec}_k(\mathcal A',\bar a')\)
for a witness \((\mathcal A',\bar a')\) of \(t\in X\); by the implication from (2)
to (3) of \cref{thm:p03-w03-p03-038-pointed-ef}, \(\varphi\) holds at
\(\bar a\). Thus \(\varphi\) and \(\bigvee_{t\in X}\chi_t\) are true at the same
tuples of all structures. Distinct equivalence classes give distinct sets
\(X\), and there are \(2^{|\mathrm{Rec}_{n,k}|}\) subsets, which is finite by
\cref{thm:p03-w03-p03-038-finite-records}. Sentences are the case \(n=0\).
\end{proof}
\end{annexitem}

\begin{annexitem}{annex:p03-w03-p03-038-ee}{thm:p03-w03-p03-038-ee}{Elementary equivalence as all finite games}
\begin{proof}
Every sentence has some finite quantifier rank. Hence
\(\mathcal A\equiv\mathcal B\) iff the structures agree through rank \(k\) for
every \(k\), iff, by \cref{thm:p03-w03-p03-038-pointed-ef} with \(n=0\),
Duplicator wins \(G_k(\mathcal A,\mathcal B)\) for every \(k\).
\end{proof}
\end{annexitem}

\begin{annexitem}{annex:p03-w03-p03-038-criterion}{thm:p03-w03-p03-038-criterion}{The game criterion}
\begin{proof}
Suppose \(\varphi\) were such a sentence and let \(k=\operatorname{qr}(\varphi)\).
Duplicator wins \(G_k(\mathcal A_k,\mathcal B_k)\), so by
\cref{thm:p03-w03-p03-038-pointed-ef} the two structures agree on all sentences
of rank at most \(k\), in particular on \(\varphi\). But
\(\mathcal A_k\in\mathcal C\) gives \(\mathcal A_k\models\varphi\), and
\(\mathcal B_k\notin\mathcal C\) gives \(\mathcal B_k\not\models\varphi\).
This contradiction proves the corollary.
\end{proof}
\end{annexitem}

\begin{annexitem}{annex:p03-w03-p03-038-cycle-response}{thm:p03-w03-p03-038-cycle-response}{Response lemma for cycles}
\begin{proof}
\emph{Cycle arithmetic.} In \(C_N\), for a vertex \(u\) and an integer \(s\)
with \(|s|\le N/2\), \(d(u,u+s)=|s|\): the vertices
\(u,u\pm1,\ldots,u+s\) give a path of length \(|s|\), and a walk of length
\(L\) from \(u\) reaches \(u+s'\) with \(|s'|\le L\) and \(s'\equiv s\pmod N\),
while every integer congruent to \(s\) modulo \(N\) has absolute value at least
\(|s|\). Consequently, if \(d(u,v)=\delta<N/2\) then \(v=u+s\) for some
\(s\) with \(|s|=\delta\), unique when \(\delta=0\) and one of the two values
\(\pm\delta\) otherwise. Vertices in different components have distance
\(\infty\). The cycle lengths here are \(2m\) and \(m\), and
\(m\ge m_k\ge2^{k+2}\ge2^{r+3}\), so every \(\delta\le2^{r+2}\) satisfies
\(\delta\le N/2\), and every \(\delta\le2^{r+1}\) satisfies \(\delta<N/2\).
The triangle inequality holds for \(d\), with \(\infty\) larger than every
integer.

\emph{Weakening.} If \(x\approx_{r+1}y\) then \(x\approx_ry\): if
\(x=y\le2^{r+1}\), either \(x=y\le2^r\) or both exceed \(2^r\); if both exceed
\(2^{r+1}\) they exceed \(2^r\). So the old pairs of the position satisfy
\(\mathrm I_r\), and the new pair has \(d(a,a)=0=d(b,b)\). It remains to choose
the answer so that \(d(a,a_j)\approx_rd(b,b_j)\) for all \(j\le p\).

We treat a move \(a\) in \(C_{2m}\); a move in \(C_m\sqcup C_m\) is handled by
the same argument with the two graphs exchanged, since \(\mathrm I_{r+1}\) is
symmetric and both graphs have all cycle lengths at least \(m\) and \(2m\)
vertices.

\emph{Close move.} Suppose \(d(a,a_i)\le2^r\) for some \(i\); fix such an
\(i\) and write \(a=a_i+t\) with \(|t|=d(a,a_i)\). Let
\(J\defeq\{j\le p:d(a_j,a_i)\le2^{r+1}\}\). For \(j\in J\) write
\(a_j=a_i+s_j\) with \(|s_j|=d(a_j,a_i)\). By \(\mathrm I_{r+1}\),
\(d(b_j,b_i)=|s_j|\le2^{r+1}\), so \(b_j\) lies in the cycle of \(b_i\) and
\(b_j=b_i+\varepsilon_js_j\) for some \(\varepsilon_j\in\{1,-1\}\), arbitrary
when \(s_j=0\).

We claim that one sign serves for all \(j\in J\) with \(s_j\ne0\). Suppose
\(j,l\in J\), \(s_j,s_l\ne0\) and \(\varepsilon_l=-\varepsilon_j\). The offsets
have absolute value at most \(2^{r+2}\), so by the cycle arithmetic
\(d(a_j,a_l)=|s_j-s_l|\) and \(d(b_j,b_l)=|\varepsilon_js_j+\varepsilon_js_l|
=|s_j+s_l|\). One of \(|s_j-s_l|,|s_j+s_l|\) equals \(|s_j|+|s_l|\) and the
other equals \(\bigl||s_j|-|s_l|\bigr|\), which is strictly smaller because
both offsets are nonzero and is at most \(2^{r+1}-1\). So one of the two
distances is at most \(2^{r+1}\) while the two differ, contrary to
\(d(a_j,a_l)\approx_{r+1}d(b_j,b_l)\). Let \(\varepsilon\) be this common sign,
or \(1\) if no \(s_j\) is nonzero; then \(b_j=b_i+\varepsilon s_j\) for all
\(j\in J\). Answer \(b\defeq b_i+\varepsilon t\).

For \(j\in J\), \(|t-s_j|\le2^r+2^{r+1}<2^{r+2}\), so
\(d(a,a_j)=|t-s_j|=|\varepsilon t-\varepsilon s_j|=d(b,b_j)\), and these equal
values satisfy \(\approx_r\) whether or not they exceed \(2^r\). For
\(j\notin J\), \(d(a_j,a_i)>2^{r+1}\), so by \(\mathrm I_{r+1}\) also
\(d(b_j,b_i)>2^{r+1}\). The triangle inequality gives
\(d(a,a_j)\ge d(a_j,a_i)-d(a,a_i)>2^{r+1}-2^r=2^r\) and likewise
\(d(b,b_j)\ge d(b_j,b_i)-d(b,b_i)>2^r\), using \(d(b,b_i)=|t|\le2^r\). Hence
\(d(a,a_j)\approx_rd(b,b_j)\).

\emph{Far move.} Suppose \(d(a,a_i)>2^r\) for all \(i\le p\) (this includes
\(p=0\)). Each ball \(\{v:d(v,b_i)\le2^r\}\) has at most \(2^{r+1}+1\le2^k+1\)
vertices, and \(p\le k-1\), so at most \((k-1)(2^k+1)\) vertices lie in some
ball. Since \(C_m\sqcup C_m\) has \(2m\ge2(k+1)2^{k+2}>(k-1)(2^k+1)\) vertices,
some \(b\) lies in none of them. Then \(d(a,a_j)>2^r\) and \(d(b,b_j)>2^r\) for
every \(j\), so \(d(a,a_j)\approx_rd(b,b_j)\). In both cases
\((\bar aa,\bar bb)\) satisfies \(\mathrm I_r\).
\end{proof}
\end{annexitem}

\begin{annexitem}{annex:p03-w03-p03-038-connectivity}{thm:p03-w03-p03-038-connectivity}{Connectivity is not first-order definable}
\begin{proof}
For \(k=0\), the empty position is a partial isomorphism and Duplicator wins.
Let \(k\ge1\) and \(m\ge m_k\). Fix enumerations of the two finite vertex sets.
Duplicator's strategy: when \(r+1\) rounds remain, answer Spoiler's move with
the first vertex in the enumeration that satisfies
\cref{thm:p03-w03-p03-038-cycle-response}. No choice principle is needed. The
empty position satisfies \(\mathrm I_k\) vacuously, and before the round with
\(r+1\) rounds remaining the position has length \(k-r-1\). By the lemma and
induction, after the last round the position of length \(k\) satisfies
\(\mathrm I_0\). Then \(d(a_i,a_j)=0\) iff \(d(b_i,b_j)=0\), since a value
\(0\le2^0\) forces equality, so \(a_i=a_j\) iff \(b_i=b_j\); and
\(d(a_i,a_j)=1\) iff \(d(b_i,b_j)=1\), so \(E(a_i,a_j)\) iff \(E(b_i,b_j)\). The
final position is a partial isomorphism, and Duplicator wins.

By \cref{thm:p03-w03-p03-038-pointed-ef}, \(C_{2m_k}\) and
\(C_{m_k}\sqcup C_{m_k}\) agree through rank \(k\). The first graph is
connected: any two vertices are joined along the cycle. The second is not:
no edge joins \((u,0)\) to \((v,1)\), so a path starting in one copy stays in
it. Suppose a sentence \(\varphi\) were true in exactly the connected finite
graphs among finite graphs. The sentence
\(\varphi\wedge\forall x\,\neg E(x,x)\wedge
\forall x\forall y\,(E(x,y)\to E(y,x))\) would be true in a finite
\(\{E\}\)-structure iff it is a connected graph. The class of connected finite
graphs, with \(\mathcal A_k=C_{2m_k}\) and
\(\mathcal B_k=C_{m_k}\sqcup C_{m_k}\), satisfies the hypothesis of
\cref{thm:p03-w03-p03-038-criterion}, which excludes such a sentence.
\end{proof}
\end{annexitem}

\begin{annexitem}{annex:p03-w03-p03-038-infinite-signature}{thm:p03-w03-p03-038-infinite-signature}{The infinite-signature counterexample}
\begin{proof}
\emph{Spoiler wins \(G_1\).} Spoiler chooses \(e\). Every answer \(i\in\N\)
satisfies \(P_i\), which \(e\) does not, so the position \((i,e)\) disagrees
on the atomic formula \(P_i(x_1)\).

\emph{Elementary equivalence.} A sentence \(\varphi\) contains finitely many
symbols, all among \(\tau_0=\{P_i:i<M\}\) for some \(M\). Call an element
\emph{blank} if it lies in no \(P_i\) with \(i<M\); the blanks of
\(\mathcal A\) are the \(i\ge M\), and those of \(\mathcal B\) are these and
\(e\); both sets are infinite. Let \(I\) be the set of finite positions between
the \(\tau_0\)-reducts that are partial isomorphisms; the empty position is in
\(I\). Let \((\bar a,\bar b)\in I\) and \(c\in\dom\mathcal A\). If
\(c=a_i\), answer \(b_i\). If \(c=i<M\) and \(c\) is new, answer \(i\): if \(i\)
were some \(b_l\), then \(P_i(b_l)\) would force \(P_i(a_l)\), so \(a_l=i=c\),
contrary to newness. If \(c\) is a new blank, answer a blank not among
\(\bar b\). In each case the extended position preserves and reflects
equality and every \(P_i\) with \(i<M\). The back condition is the same
argument with the structures exchanged, using that \(\mathcal A\) also has
infinitely many blanks. So \(I\) is a back-and-forth family, Duplicator wins
every \(G_k\) between the reducts by \cref{thm:p03-w03-p03-038-bf-system}, and
the reducts agree on every \(\tau_0\)-sentence by
\cref{thm:p03-w03-p03-038-ee}, the signature \(\tau_0\) being finite. As
\(\varphi\) is a \(\tau_0\)-sentence, its truth value in each structure equals
that in its reduct. Hence \(\mathcal A\equiv\mathcal B\).
\end{proof}
\end{annexitem}

\begin{annexitem}{annex:p03-w03-p03-038-bf-system}{thm:p03-w03-p03-038-bf-system}{Systems, games and records}
\begin{proof}
\emph{(1).} Suppose Duplicator wins
\(G_k(\mathcal A,\bar a;\mathcal B,\bar b)\). Let \(I_j\) be the set of
positions from which Duplicator has a winning strategy in the \(j\)-round
game. Members of \(I_0\) are partial isomorphisms by definition; members of
\(I_j\) with \(j>0\) are restrictions of final positions of won plays, hence
partial isomorphisms. If \(P\in I_{j+1}\) is won by \(\sigma\) and Spoiler
moves \(c\) or \(d\), \(\sigma\)'s answer leads to a position won by the
continuation of \(\sigma\) in \(j\) rounds, which therefore lies in \(I_j\).
So \((I_j)_{j\le k}\) is a system with \((\bar a,\bar b)\in I_k\). Conversely,
given a system with \((\bar a,\bar b)\in I_k\), Duplicator answers, while
\(j+1\) rounds remain at a position in \(I_{j+1}\), with an answer leading
into \(I_j\), fixed by the Axiom of Choice. The final position lies in
\(I_0\) and is a partial isomorphism.

\emph{(2).} Equal \(j\)-records have equal first components, so each member of
\(I_j\) is a partial isomorphism. If \(P=(\bar a',\bar b')\in I_{j+1}\), the
second components of the \((j+1)\)-records coincide, so for each \(c\) some
\(d\) has \(\operatorname{rec}_j(\mathcal A,\bar a'c)=\operatorname{rec}_j(\mathcal B,\bar b'd)\),
and conversely; this is forth and back. If \((I'_j)_{j\le k}\) is any system
and \(P\in I'_j\), then \((I'_i)_{i\le j}\) is a \(j\)-system containing
\(P\) in its top level, so Duplicator wins \(G_j\) from \(P\) by (1), and
\(P\in I_j\) by \cref{thm:p03-w03-p03-038-pointed-ef}.

\emph{(3).} Combine (1) for the empty position with
\cref{thm:p03-w03-p03-038-ee}.
\end{proof}
\end{annexitem}

\begin{annexitem}{annex:p03-w03-p03-038-fraisse-bridge}{thm:p03-w03-p03-038-fraisse-bridge}{Fra\"iss\'e's extension condition gives the comparison games}
\begin{proof}
Read \(R\) and \(S\) as structures for one \(n\)-ary symbol. An embedding in
the sense of \cref{def:p03-fraisse54-relation} is an injection preserving
the value of every \(n\)-tuple, so a finite position is a partial
isomorphism iff \(a_i\mapsto b_i\) is a well-defined embedding of
\(R\mathbin{\upharpoonright}\{a_1,\ldots,a_m\}\) into \(S\).

Let \(F\) be the set of finite positions that are partial isomorphisms; it
contains the empty position. Let \((\bar a,\bar b)\in F\) and \(c\) be in the
base of \(R\). If \(c=a_i\), the answer \(b_i\) keeps the position in \(F\),
since every atom mentioning the new coordinate is equivalent to one with
coordinate \(i\) in its place. Otherwise put \(A_0\defeq\{a_1,\ldots,a_m\}\),
\(A\defeq R\mathbin{\upharpoonright}A_0\) and
\(B\defeq R\mathbin{\upharpoonright}(A_0\cup\{c\})\). Both are finite
restrictions of \(R\), hence members of \(\gamma_R=\gamma_S\), and \(A\) is a
restriction of \(B\). The map \(\varphi:a_i\mapsto b_i\) is an embedding of
\(A\) into \(S\). Condition \(\mathcal H\) for \(S\)
(\cref{def:p03-fraisse54-extension}) extends \(\varphi\) to an embedding
\(\psi\) of \(B\) into \(S\). Answer \(d\defeq\psi(c)\); the position
\((\bar ac,\bar bd)\) corresponds to \(\psi\) and lies in \(F\). The back
condition is the same argument with \(R\) and \(S\) exchanged, using
\(\mathcal H\) for \(R\). Thus \(F\) is a back-and-forth family; no
cardinality hypothesis was used.

Taking \(I_j=F\) for all \(j\) gives a \(k\)-back-and-forth system for every
\(k\), so by \cref{thm:p03-w03-p03-038-bf-system} Duplicator wins
\(G_k(R,\bar a;S,\bar b)\) for each position in \(F\); for the empty
position, \cref{thm:p03-w03-p03-038-ee} gives \(R\equiv S\).

If both bases are finite or countable, list them as \(d_0,d_1,\ldots\) and
\(e_0,e_1,\ldots\), repeating entries if a base is finite. Starting from a
position in \(F\), Spoiler chooses \(d_0,e_0,d_1,e_1,\ldots\), and Duplicator
answers with the first element of the other list that keeps the position in
\(F\); the forth and back conditions guarantee such an element, and no
choice principle is needed. Every finite subposition of the union lies in
\(F\), so the union is a partial isomorphism whose domain is the whole base
of \(R\) and whose range is the whole base of \(S\). It is therefore a
bijection preserving every value, an isomorphism of \(R\) onto \(S\) that
extends the initial position. This is the conclusion of
\S5.4(2\textsuperscript{o}) of the source.
\end{proof}
\end{annexitem}

\begin{annexitem}{annex:p03-w03-p03-038-e61-determined}{thm:p03-w03-p03-038-e61-determined}{The games are determined}
\begin{proof}
We prove by induction on the number \(j\) of remaining moves that from every
position exactly one player has a winning strategy; both players cannot,
since the two strategies played against each other would produce a play won
by both. For \(j=0\) the play is over, and II wins iff the correspondence is
an isomorphism for \(A\); the winner needs no choices. Let \(j+1\) moves
remain at a position \(P\). Suppose that for every legal move \(m\) of I
(a model and a sequence, or a single individual in \(G_n\)) some answer of
II leads to a position from which II wins in \(j\) moves. By the Axiom of
Choice fix one such answer for each \(m\) and a winning strategy from each
resulting position; together they form a winning strategy for II from
\(P\). Otherwise some move \(m\) has the property that every answer leads to
a position from which, by the induction hypothesis, I has a winning
strategy; choosing one for each answer, I plays \(m\) and then follows the
chosen strategy, and wins.
\end{proof}
\end{annexitem}

\begin{annexitem}{annex:p03-w03-p03-038-e61-int-rat}{thm:p03-w03-p03-038-e61-int-rat}{The strategies of the example}
\begin{proof}
\emph{II wins \(H_1\).} Let I choose \(a_1,\ldots,a_k\) in one model. Let
\(c_0<\cdots<c_{p-1}\) be the distinct values among them, and let II answer
\(a_j\) by the integer \(i\), which is also a rational, where \(a_j=c_i\).
Then \(a_j=a_{j'}\) iff the answers are equal, and \(a_j<a_{j'}\) iff the
answers satisfy \(<\); \(\eta\) holds of all individuals. So the
correspondence is an isomorphism for \(\eta,\iota,\varrho\).

\emph{I wins \(G_3\).} I chooses the integer \(0\), then the integer \(1\);
let \(w_1,w_2\) be II's rational answers. If \(w_1<w_2\) fails, the pairs
\(0\leftrightarrow w_1\), \(1\leftrightarrow w_2\) already violate
\(\varrho\), whatever happens next. If \(w_1<w_2\), I chooses the rational
\((w_1+w_2)/2\), which lies strictly between them. An integer answer \(x\)
would need \(0<x<1\) to keep \(\varrho\), and no integer satisfies this.
\end{proof}
\end{annexitem}

\begin{annexitem}{annex:p03-w03-p03-038-e61-hull}{thm:p03-w03-p03-038-e61-hull}{A matching of individuals extends to the finite sets over them}
\begin{proof}
Write \(I_1\) for the individuals of \(\mathfrak M_1\). By clause~(2),
\(|\mathfrak M_1|=\bigcup_rV_r\) with \(V_0=I_1\) and \(V_{r+1}\) the union
of \(V_r\) with the finite nonempty subsets of \(V_r\); similarly
\(\operatorname{HF}(X)=\bigcup_rV_r(X)\) with \(V_0(X)=X\). Call the least
\(r\) with \(x\in V_r\) the \emph{stage} of \(x\). By the reading
\cref{def:p03-w03-p03-038-e61-reading}(1), an element of stage \(r+1\) is a
non-individual all of whose members have stage at most \(r\), and an
individual has no members.

\emph{The description of \(\operatorname{HF}(X)\).} By induction on stage,
every \(x\) has finite \(S(x)\) and \(s(x)\): \(S(x)=\{x\}\) for an
individual, and \(S(x)=\{x\}\cup\bigcup_{t\in x}S(t)\) otherwise, a finite
union of finite sets, because the right side contains \(x\), is closed
under members and is contained in every such set. If \(x\in V_r(X)\), then
\(s(x)\subseteq X\), by induction on \(r\). Conversely, if
\(s(x)\subseteq X\), then \(x\in\operatorname{HF}(X)\) by induction on the
stage of \(x\): an individual \(x\) lies in \(s(x)\subseteq X\), and a
non-individual \(x\) has members \(t\) with \(s(t)\subseteq s(x)\subseteq
X\), hence \(t\in\operatorname{HF}(X)\), and \(x\) is a finite nonempty set
of them.

\emph{The map.} The correspondence respects \(\iota\in A\), so
\(f:a_i\mapsto b_i\) is a well-defined bijection of \(X\) onto \(Y\).
Define \(F\) on \(V_r(X)\) by recursion on \(r\): \(F(a_i)=b_i\), and
\(F(x)=\{F(t):t\in x\}\) for a non-individual \(x\). The two clauses do not
overlap, so \(F\) is well defined and unique, and \(F(x)\in V_r(Y)\) when
\(x\in V_r(X)\). Individuals go to individuals, non-individuals to
non-individuals. \(F\) is injective: if \(F(x)=F(x')\), both are individuals
and \(f\) is injective, or both are sets, and then each \(t\in x\) has
\(t'\in x'\) with \(F(t)=F(t')\), so \(t=t'\) by induction on stage; hence
\(x\subseteq x'\), and symmetrically. \(F\) is onto \(\operatorname{HF}(Y)\)
by induction on the stage in \(\operatorname{HF}(Y)\), using \(f^{-1}\) for
individuals.

\emph{Atomic formulas.} For \(\iota\), use injectivity. For \(\epsilon\):
\(t\in x\) gives \(F(t)\in F(x)\) by definition; if \(F(t)\in F(x)\), then
\(F(x)\) is not an individual, so \(x\) is a set and \(F(t)=F(t')\) for
some \(t'\in x\), whence \(t=t'\in x\). For \(\alpha\in A\): by clause~(3)
\(\alpha\) holds only of individuals, and \(F\) preserves being an
individual, so if some argument is not an individual both sides fail; if
all arguments are individuals \(a_{i_1},\ldots\), then \(\alpha\) holds of
them iff it holds of \(b_{i_1},\ldots\), because the correspondence is an
isomorphism for \(A\).
\end{proof}
\end{annexitem}

\begin{annexitem}{annex:p03-w03-p03-038-e61-thm1}{thm:p03-w03-p03-038-e61-thm1}{Theorems 1--3 of the source}
\begin{proof}
\emph{(1).} Induction on \(n-k\), following the source. If \(n-k=0\),
\(\Psi\) is quantifier-free. Let \(F\) be the map of
\cref{thm:p03-w03-p03-038-e61-hull} for the correspondence of the \(k\)
moves; it sends \(x_i\) to \(y_i\), and it preserves every atomic formula,
so it preserves every Boolean combination of them, which gives the claim.

Let \(n-k-1\ge0\) and assume the claim for \(k+1\) free variables. A
formula \((\forall u)\Phi\) is equivalent to \(\sim(\exists u)\Phi'\) where
\(\Phi'\in\mathfrak H_n(A)\) is obtained from \(\sim\Phi\) by exchanging
every prefix quantifier with its dual and negating the matrix; it has the
same free variables and the same number of quantifiers. So it suffices to
treat \(\Psi=(\exists u)\Phi(x_1,\ldots,x_k,u)\) with
\(\Phi\in\mathfrak H_n(A)\) having \(k+1\) free variables, and to show that
truth in \(\mathfrak M_1\) gives truth in \(\mathfrak M_2\); the other
direction is the same argument with the models exchanged. Let
\(x_{k+1}\) witness \(\Phi\) in \(\mathfrak M_1\). In move \(k+1\), I
chooses \(\mathfrak M_1\) and lists \(s(x_{k+1})\), or any single individual
if \(s(x_{k+1})\) is empty (individuals exist by
\cref{def:p03-w03-p03-038-e61-reading}(2)). II answers by her winning
strategy. The correspondence of the \(k+1\) moves is an isomorphism for
\(A\): it is the restriction of the final correspondence of a play
continuing it under her strategy, and every finite continuation is legal.
Its map \(F'\) extends the map \(F\) of the first \(k\) moves, by the
uniqueness in \cref{thm:p03-w03-p03-038-e61-hull}, so it still sends
\(x_i\) to \(y_i\); and \(s(x_{k+1})\) lies in its domain. Put
\(y_{k+1}\defeq F'(x_{k+1})\). Then the correspondence establishes an
isomorphism between \(x_i\) and \(y_i\) for \(i\le k+1\), II has a winning
strategy after these \(k+1\) moves, and the induction hypothesis gives
\(\operatorname{stsf}_{\mathfrak M_2}\Phi(y_1,\ldots,y_{k+1})\), hence
\(\operatorname{stsf}_{\mathfrak M_2}\Psi(y_1,\ldots,y_k)\).

\emph{(2).} Take \(k=0\) in (1); the empty correspondence establishes
nothing and is needed for nothing.

\emph{(3).} \emph{Prenex forms.} In structures with nonempty domains
every formula is equivalent to a prenex formula with the same free
variables whose prefix quantifiers bind distinct variables different from
the free ones. By induction on formulas: atomic formulas and \(\bottom\)
are prenex. For \(\sim\varphi\), take a prenex form of \(\varphi\) and use
\(\sim(\exists u)\theta\equiv(\forall u)\sim\theta\) and
\(\sim(\forall u)\theta\equiv(\exists u)\sim\theta\) along the prefix. For a
binary connective, take prenex forms of both sides, rename bound variables
so that no variable bound on one side occurs on the other, rewrite
\(\equiv\) as a conjunction of two implications, and move the quantifiers
out one at a time by the equivalences, valid when \(u\) is not free in
\(\psi\) and the domain is nonempty,
\((Qu)\theta\wedge\psi\equiv(Qu)(\theta\wedge\psi)\),
\((Qu)\theta\vee\psi\equiv(Qu)(\theta\vee\psi)\),
\(\psi\supset(Qu)\theta\equiv(Qu)(\psi\supset\theta)\) and
\((Qu)\theta\supset\psi\equiv(Q^*u)(\theta\supset\psi)\), where \(Q^*\) is
the dual quantifier, together with their mirror images. For
\((Qu)\varphi\), take a prenex form of \(\varphi\) whose bound variables
differ from \(u\) and prefix \(Qu\). Each satisfaction clause used is
checked directly, nonemptiness being needed only to move a quantifier past
a formula in which its variable is not free.

A closed formula of \(\mathfrak H(A)\) is therefore equivalent, in every
model, to a closed prenex formula, which with \(n\) quantifiers lies in
\(\mathfrak H_n(A)\); and (2) applies.
\end{proof}
\end{annexitem}

\begin{annexitem}{annex:p03-w03-p03-038-e61-thm4}{thm:p03-w03-p03-038-e61-thm4}{Theorem 4 of the source}
\begin{proof}
\(\mathfrak H(A)\) has countably many formulas. For each formula
\(\varphi(u,v_1,\ldots,v_p)\) fix by the Axiom of Choice a function
\(f_\varphi:|\mathfrak M|^p\to|\mathfrak M|\) with
\(\operatorname{stsf}_{\mathfrak M}\varphi(f_\varphi(\bar c),\bar c)\)
whenever \(\operatorname{stsf}_{\mathfrak M}(\exists u)\varphi(u,\bar c)\).
Fix an individual \(x_0\) and let \(X\) be the least set containing \(x_0\)
and closed under (a) forming \(\{z_1,\ldots,z_k\}\) from
\(z_1,\ldots,z_k\in X\), (b) passing from \(z\in X\) to its members, and
(c) the functions \(f_\varphi\). Then \(X=\bigcup_rX_r\), with
\(X_0=\{x_0\}\) and \(X_{r+1}\) obtained from \(X_r\) by applying (a)--(c)
once; each \(X_{r+1}\) is countable, being a countable union of countable
sets (countable choice), so \(X\) is countable. It is infinite: \(x_0\),
\(\{x_0\}\), \(\{\{x_0\}\},\ldots\) are distinct by the Axiom of
Foundation and the reading that individuals are not sets of clause~(2).

Let \(\mathfrak M'\) have domain \(X\), with every predicate and \(\epsilon\)
restricted to \(X\). It is an elementary substructure of \(\mathfrak M\): by induction on
formulas, a formula with parameters from \(X\) holds in \(\mathfrak M'\)
iff it holds in \(\mathfrak M\). Atomic formulas are evaluated by the
restricted interpretations, and the connectives are immediate. If
\((\exists u)\varphi(u,\bar c)\) holds in \(\mathfrak M'\), a witness in
\(X\) satisfies \(\varphi\) in \(\mathfrak M\) by the induction hypothesis;
if it holds in \(\mathfrak M\), then \(f_\varphi(\bar c)\in X\) by (c)
satisfies \(\varphi\) in \(\mathfrak M\), hence in \(\mathfrak M'\). The
universal quantifier is the dual case. It is a model for \(\mathfrak H(A)\): its
individuals are \(X\cap\hat x\eta(x)\), the predicates of \(A\) live on
them, \(\iota\) is identity and \(\epsilon\) is membership. For clause~(2),
\(\operatorname{HF}(X\cap\hat x\eta(x))\subseteq X\) by (a), and conversely
each \(z\in X\) has \(S(z)\subseteq X\) by (b), so
\(s(z)\subseteq X\cap\hat x\eta(x)\) and
\(z\in\operatorname{HF}(X\cap\hat x\eta(x))\) by
\cref{thm:p03-w03-p03-038-e61-hull}.
\end{proof}
\end{annexitem}

\begin{annexitem}{annex:p03-w03-p03-038-e61-sec3}{thm:p03-w03-p03-038-e61-sec3}{Indiscernible models on which I wins \(H_3\)}
\begin{proof}
\(\mathfrak M_2\) is an elementary submodel, so the two models satisfy the
same closed formulas of \(\mathfrak H(A)\).

\emph{Every natural number lies in \(\mathfrak M_2\).} For each \(m\), the
formula \(\varphi_m(x)\): ``\(\beta(x)\) and exactly \(m\) individuals \(y\)
satisfy \(\gamma(y,x)\)'' is first-order and holds in \(\mathfrak M_1\)
only of \(m\). As \(\mathfrak M_1\) satisfies \((\exists x)\varphi_m(x)\),
so does \(\mathfrak M_2\), and a witness \(c\) in \(\mathfrak M_2\)
satisfies \(\varphi_m\) in \(\mathfrak M_1\) by elementarity; so \(c=m\).
The individuals of \(\mathfrak M_1\) include uncountably many sequences
(Cantor's diagonal argument) and \(\mathfrak M_2\) is countable, so some
sequence \(x\) of \(\mathfrak M_1\) is not in \(\mathfrak M_2\).

\emph{I's strategy.} First move: \(x\) in \(\mathfrak M_1\); let \(y\) be
the answer. If \(\alpha(y)\) fails, the correspondence violates \(\alpha\)
and I has won. Otherwise \(y\) is a sequence different from \(x\); choose
\(m\) with \(y_m\ne x_m\). Then exactly one of \(\delta(x,m,0)\) in
\(\mathfrak M_1\) and \(\delta(y,m,0)\) in \(\mathfrak M_2\) holds. Second
move: \(0,1,\ldots,m\) in \(\mathfrak M_1\); let \(a_0,\ldots,a_m\) be the
answers. (a) If \(a_i=i\) for all \(i\), the pairs \(x\leftrightarrow y\),
\(m\leftrightarrow m\), \(0\leftrightarrow0\) violate \(\delta\).
(b) If some \(a_i\) is not a natural number, they violate \(\beta\).
(c) If all are natural numbers and \(a_i<a_j\) fails for some \(i<j\),
they violate \(\gamma\). (d) Otherwise \(a_0<\cdots<a_m\) are natural
numbers, not all \(a_i=i\). Then \(a_i\ge i\) for all \(i\), and if
\(a_i>i\) then \(a_j>j\) for \(j\ge i\); so \(a_m>m\), and among the
\(a_m\ge m+1\) numbers below \(a_m\) at most the \(m\) numbers
\(a_0,\ldots,a_{m-1}\) are chosen. Third move: such an unchosen \(b<a_m\),
in \(\mathfrak M_2\). An answer \(c\) in \(\mathfrak M_1\) keeping the
correspondence an isomorphism would satisfy \(\gamma(c,m)\) and \(c\ne i\)
for \(i\le m\), which no individual does. In cases (a)--(c) the third move
is arbitrary. In every case the final correspondence is not an isomorphism.
\end{proof}
\end{annexitem}

\begin{annexitem}{annex:p03-w03-p03-038-e61-thm5}{thm:p03-w03-p03-038-e61-thm5}{Theorems 5--7 of the source}
\begin{proof}
\emph{(1).} Induction on \(n-k\). The \(x_i,y_i\) are individuals and the
correspondence of the \(k\) moves is an isomorphism for \(A\), as in the
proof of \cref{thm:p03-w03-p03-038-e61-thm1}. A quantifier-free formula of
\(\mathfrak G(A)\) contains no \(\epsilon\), so its atoms are formulas
\(\alpha(\ldots)\) with \(\alpha\in A\), on which the correspondence agrees.
For the step, by duality as before let \(\Psi=(\exists u)_\eta\Phi\), and
let the individual \(x_{k+1}\) witness \(\Phi\) in \(\mathfrak M_1\). I
chooses \(x_{k+1}\); II answers \(y_{k+1}\) by her strategy; she has a
winning strategy after \(k+1\) moves, and the induction hypothesis gives
\(\Phi(y_1,\ldots,y_{k+1})\) in \(\mathfrak M_2\), so \(\Psi\) holds there.
The other direction exchanges the models.

\emph{(2).} Take \(k=0\).

\emph{(3).} A formula of \(\mathfrak G(A)\) is evaluated in \(\mathfrak M\)
exactly as the formula obtained by deleting the restrictions is evaluated in
the structure \(\mathfrak M^\eta\) of individuals with the relations of
\(A\), by induction on formulas: a restricted quantifier ranges over
individuals. By the prenex forms of the proof of
\cref{thm:p03-w03-p03-038-e61-thm1}(3), a closed formula is equivalent in
\(\mathfrak M^\eta\), whose domain is nonempty, to a closed prenex formula;
restoring the restrictions gives a closed formula of some
\(\mathfrak G_n(A)\) with the same truth value in \(\mathfrak M\). Apply
(2).
\end{proof}
\end{annexitem}

\begin{annexitem}{annex:p03-w03-p03-038-unary-types}{thm:p03-w03-p03-038-unary-types}{Counting unary types}
\begin{proof}
A finite position is a partial isomorphism iff it respects equality and
pairs elements of the same type \(t\). Let \((\bar a,\bar b)\) be one and
\(c\) an element of \(\mathcal A\) of type \(t\). If \(c=a_i\), answer
\(b_i\). Otherwise the position pairs the chosen elements of type \(t\)
bijectively, so the same number \(q\) of them is chosen on each side, and
\(n_t(\mathcal A)>q\). Then \(n_t(\mathcal B)>q\) as well, because the two
counts are equal or both infinite, and an unchosen element of
\(\mathcal B\) of type \(t\) is a correct answer. The back condition is
symmetric. The family contains the empty position, so
\cref{thm:p03-w03-p03-038-bf-system} gives a winning strategy in every
\(G_k\), and \cref{thm:p03-w03-p03-038-ee} gives
\(\mathcal A\equiv\mathcal B\), the signature \(P\) being finite.
\end{proof}
\end{annexitem}

\begin{annexitem}{annex:p03-w03-p03-038-e61-converse-fails}{thm:p03-w03-p03-038-e61-converse-fails}{The converse of Theorem 7 fails for infinite \(A\)}
\begin{proof}
\emph{I wins \(G_1\).} He chooses \(\omega\) in \(\mathfrak M_2\). Any
answer \(m\) satisfies \(\sim a_m(m)\) (if \(m\ge1\)) while \(a_m(\omega)\),
or, if \(m=0\), \(\sim a_1(0)\) while \(a_1(\omega)\).

\emph{Elementary indiscernibility.} A closed formula of \(\mathfrak G(A)\)
mentions only \(a_1,\ldots,a_M\) for some \(M\). Its truth value in each
model is that of the unrestricted formula in the structure of individuals
with \(a_1,\ldots,a_M\) (proof of \cref{thm:p03-w03-p03-038-e61-thm5}(3)).
In these structures, the type \(\{a_i:i<m\}\) is realized by \(m\) alone for
\(2\le m\le M\), the empty type by \(0\) and \(1\), the full type
\(\{a_1,\ldots,a_M\}\) by the infinitely many \(m>M\) and, in
\(\mathfrak M_2\), also by \(\omega\); no other type is realized. The counts
agree or are both infinite, so \cref{thm:p03-w03-p03-038-unary-types} gives
equivalence for \(\{a_1,\ldots,a_M\}\).
\end{proof}
\end{annexitem}

\begin{annexitem}{annex:p03-w03-p03-038-discrete-orders}{thm:p03-w03-p03-038-discrete-orders}{Gap strategy for discrete orders}
\begin{proof}
For a finite set \(c_1<\cdots<c_q\) of selected points of \(L\), the gaps
are \(\{x:x<c_1\}\), \(\{x:c_i<x<c_{i+1}\}\) and \(\{x:x>c_q\}\) (the whole
of \(L\) if \(q=0\)), with sizes in \(\N\cup\{\infty\}\). Write
\(g\simeq_rg'\) if \(g=g'\) or both are at least \(2^r-1\). A position
satisfies \(J_r\) if it is a partial isomorphism for \(\{<\}\), so that
its two sets of selected points correspond in order, and corresponding
gaps have \(\simeq_r\) sizes. The empty position satisfies \(J_k\): both
orders are infinite. \(J_0\) implies being a partial isomorphism.

\emph{Points at the ends of a gap.} Let \(G\) be a gap of \(L\) with at
least \(p\) points. Its \(p\) least points exist: the least is the least
element of \(L\) if \(G\) is the bottom gap, and otherwise the immediate
successor of its lower boundary; each further one is the immediate
successor of the previous one, which lies in \(G\) while fewer than \(p\)
have been listed. If \(G\) has an upper boundary, its \(p\) greatest points
exist likewise, using immediate predecessors. The top gap is infinite above
every point because \(L\) has no greatest element.

\emph{Response.} Let \(r\ge1\), let the position satisfy \(J_r\), and let
Spoiler choose \(a\) in \(L\); the other case is symmetric. If \(a\) is
selected, answer its partner; the gaps do not change and \(\simeq_r\)
implies \(\simeq_{r-1}\). Otherwise \(a\) lies in a gap \(G\) of size
\(g\), with \(\ell\) points of \(G\) below \(a\) and \(\rho\) above, and
the corresponding gap \(G'\) has size \(g'\simeq_rg\). If \(g=g'\) is
finite, answer the point of \(G'\) with exactly \(\ell\) points of \(G'\)
below it. If \(g,g'\ge2^r-1\): when \(\ell<2^{r-1}-1\), answer the
\((\ell+1)\)th least point of \(G'\), leaving \(\ell\) points below and at
least \(2^r-1-\ell-1\ge2^{r-1}-1\) above, while \(a\) has
\(\rho\ge2^{r-1}-1\) points above for the same reason; when
\(\rho<2^{r-1}-1\), then \(G\) is not the top gap and neither is \(G'\),
and answer the \((\rho+1)\)th greatest point of \(G'\); otherwise answer the
\(2^{r-1}\)th least point of \(G'\), which has \(2^{r-1}-1\) points of
\(G'\) below and at least \(2^r-1-2^{r-1}=2^{r-1}-1\) above. In every case
the answer lies in \(G'\), so the order relations with the selected points
are preserved, the two new gaps are \(\simeq_{r-1}\) on the two sides, and
the other gaps are unchanged. So \(J_{r-1}\) holds.

By induction on the rounds, after \(k\) rounds \(J_0\) holds and Duplicator
has won; she needs no choice principle, the answers being determined. By
\cref{thm:p03-w03-p03-038-ee} the orders are elementarily equivalent.
\end{proof}
\end{annexitem}

\begin{annexitem}{annex:p03-w03-p03-038-e61-thm8}{thm:p03-w03-p03-038-e61-thm8}{Theorems 8--10 of the source}
\begin{proof}
\emph{Properties of the lists.} For level \(n\) they are assumed, and the
proposed list has them: the sign patterns over all atomic formulas in
\(u_1,\ldots,u_n\), finitely many as \(A\) is finite, are pairwise
exclusive, their disjunction is a tautology, and each fixes the truth value
of every atom and hence of every Boolean combination \(\Gamma\). For level
\(k-1\), the \(\Omega_{k-1,l}\) are all sign patterns over the formulas
\((\exists u_k)\Omega_{k0},\ldots,(\exists u_k)\Omega_{kl_k}\); so they are
pairwise exclusive and their disjunction is a tautology. Hence in every
model each tuple of individuals satisfies exactly one \(\Omega_{ki}\).

\emph{(1).} Induction on \(n-k\). For \(k=n\), II wins after \(n\) moves iff
the correspondence is an isomorphism, iff \(\bar x\) and \(\bar y\)
satisfy the same atomic formulas. If they do, \(\bar y\) satisfies
\(\Omega_{ni}\). If \(\bar y\) satisfies \(\Omega_{ni}\), then every atom
\(\Gamma\) is implied, or its negation is implied, by \(\Omega_{ni}\), so
\(\bar x\) and \(\bar y\) agree on it.

Let \(1\le k\le n\), assume (1) for \(k\), and let
\(\bar x=(x_1,\ldots,x_{k-1})\) satisfy \(\Omega_{k-1,i_0}\). Suppose
\(\bar y\) satisfies \(\Omega_{k-1,i_0}\). If I chooses \(x\) in
\(\mathfrak M_1\), let \(\Omega_{ki_1}\) be the formula satisfied by
\(\bar xx\). Then \((\exists u_k)\Omega_{ki_1}\) holds at \(\bar x\), so it
is a positive conjunct of \(\Omega_{k-1,i_0}\) and holds at \(\bar y\); II
answers some \(y\) with \(\Omega_{ki_1}(\bar y,y)\), and by the induction
hypothesis she has a winning strategy after these \(k\) moves. If I chooses
\(y\) in \(\mathfrak M_2\), let \(\Omega_{ki_1}\) be satisfied by
\(\bar yy\); the same argument with the models exchanged gives \(x\) with
\(\Omega_{ki_1}(\bar x,x)\), and the induction hypothesis applies since
\(\bar xx\) satisfies \(\Omega_{ki_1}\). Fixing these answers and
strategies by the Axiom of Choice gives II a winning strategy after
\(k-1\) moves.

Suppose \(\bar y\) does not satisfy \(\Omega_{k-1,i_0}\). Then some
\((\exists u_k)\Omega_{ki_1}\) holds at exactly one of \(\bar x,\bar y\).
If it holds at \(\bar x\), I chooses \(x\) with \(\Omega_{ki_1}(\bar x,x)\);
no answer \(y\) has \(\Omega_{ki_1}(\bar y,y)\), so by the induction
hypothesis II has no winning strategy after the \(k\) moves. If it holds at
\(\bar y\), I chooses \(y\) with \(\Omega_{ki_1}(\bar y,y)\); every answer
\(x\) has \(\bar xx\) satisfying some \(\Omega_{ki_2}\) with
\(i_2\ne i_1\), and \(\bar yy\) does not satisfy \(\Omega_{ki_2}\), so by
the induction hypothesis applied to \(\Omega_{ki_2}\) II has no winning
strategy after the \(k\) moves. By
\cref{thm:p03-w03-p03-038-e61-determined}, I then has a winning strategy
from each such position, and combining them with his choice of move
(Axiom of Choice) gives him a winning strategy after \(k-1\) moves; so II
has none.

\emph{(2).} Let \(\Phi\) be the \(\Omega_{ki}\) satisfied by \(\bar x\) and
apply (1).

\emph{(3).} Let \(\Phi_1,\ldots,\Phi_s\) be the closed formulas
\(\Omega_{0,0},\ldots,\Omega_{0,l_0}\), and let \(\mathfrak M_1\) satisfy
\(\Omega_{0i}\). By (1) with \(k=0\), II wins iff \(\mathfrak M_2\)
satisfies \(\Omega_{0i}\), iff (each model satisfying exactly one of the
list) both models satisfy the same members of the list.

\emph{(4).} If the models are elementarily indiscernible they agree on
every \(\Omega_{0i}\), which are closed formulas of \(\mathfrak G(A)\), and
(3) gives a winning strategy in every \(G_n\). The converse is
\cref{thm:p03-w03-p03-038-e61-thm5}(3).
\end{proof}
\end{annexitem}

\begin{annexitem}{annex:p03-w03-p03-038-e61-remark}{thm:p03-w03-p03-038-e61-remark}{The source's remark after Theorem 7}
\begin{proof}
\emph{(1).} The models of \cref{thm:p03-w03-p03-038-e61-sec3} satisfy the
same closed formulas of \(\mathfrak H(A)\), hence of
\(\mathfrak G(A)\subseteq\mathfrak H(A)\). Their \(A\) is finite, so II
wins every \(G_n\) by \cref{thm:p03-w03-p03-038-e61-thm8}(4), while I wins
\(H_3\).

\emph{(2).} The orders \(\omega\) and \(\omega+\omega^*+\omega\) satisfy
\((\mathrm D)\): each has a least element and no greatest; in the middle
copy of \(\omega^*\) every element has an immediate successor and
predecessor, the greatest element of that copy is the immediate
predecessor of the least element of the last copy, and every non-least
element of the first and last copies has an immediate predecessor. By
\cref{thm:p03-w03-p03-038-discrete-orders} and the evaluation of
\(\mathfrak G(A)\) in the structure of individuals
(\cref{thm:p03-w03-p03-038-e61-thm5}(3)), the models are elementarily
indiscernible. The closed formula
\[
 (\forall x)\bigl(\eta(x)\supset(\exists s)(\forall y)\bigl(\eta(y)\supset
 (\epsilon(y,s)\equiv\varrho(y,x))\bigr)\bigr)
\]
of \(\mathfrak H(A)\) holds in the first model, where \(s\) is the finite
set of predecessors of \(x\) (or any \(s\) without individual members when
\(x\) is least, for instance \(\{\{x\}\}\)). It fails in the second at the
greatest element \(x\) of the copy of \(\omega^*\), which has infinitely
many predecessors, while every element of the model has finitely many
members.
\end{proof}
\end{annexitem}

\begin{annexitem}{annex:p03-w03-p03-038-e61-wins}{thm:p03-w03-p03-038-e61-wins}{Properties of \(W_j\)}
\begin{proof}
\emph{(1).} First, \(W_j(P\cup Q)\) implies \(W_j(P)\), by induction on
\(j\). For \(j=0\): the condition for an isomorphism for \(A\) quantifies
over choices of pairs, and every choice from \(P\) is a choice from
\(P\cup Q\). For \(j+1\): given a move of I, \(W_{j+1}(P\cup Q)\) yields
an answer \(R\) with \(W_j(P\cup Q\cup R)\), and the induction hypothesis,
applied to \(P\cup R\) and \(Q\), gives \(W_j(P\cup R)\); so the same
answer serves from \(P\). Second, let \(W_{j+1}(P)\). If one of the
structures has an individual, I may choose it; the answer \(R\) given by
\(W_{j+1}(P)\) satisfies \(W_j(P\cup R)\), hence \(W_j(P)\) by the first
part. If neither has individuals, \(P=\emptyset\), the empty position is an
isomorphism, and no move of I exists, so \(W_j(\emptyset)\) holds for every
\(j\). Third, applying the second part \(j\) times turns \(W_j(P)\) into
\(W_0(P)\).

\emph{(2).} By induction on \(j\): \(\pi_1,\pi_2\) map the isomorphisms for
\(A\) between the first pair onto those between the second, and map moves
and answers of each game onto moves and answers of the other.

\emph{(3).} Let \(j=n-k\). Suppose II has a winning strategy \(\tau\) from
\(P\). By induction on \(j\), \(W_j(P)\): for \(j=0\) the play is over and
\(\tau\) wins, so \(P\) is an isomorphism; for \(j+1\), every move of I is
answered by \(\tau\) with a position from which \(\tau\) still wins in
\(j\) moves, so \(W_j\) holds there. Conversely, suppose \(W_j(P)\). By the
Axiom of Choice fix, for every position \(Q\) and \(i\) with
\(W_{i+1}(Q)\) and every move of I, an answer after which \(W_i\) holds.
The strategy that uses these answers keeps \(W_{j-r}\) true after \(r\)
further moves, so after \(j\) moves \(W_0\) holds, and II wins.

\emph{(4).} Induction on \(j\), following the proof of
\cref{thm:p03-w03-p03-038-e61-thm1}(1) with \(W\) in place of the
strategy. If \(\Psi\) has no prefix quantifier, \(W_j(P)\) makes \(P\)
an isomorphism for \(A\) by (1); the map \(F\) of
\cref{thm:p03-w03-p03-038-e61-hull} sends \(x_i\) to \(y_i\) and preserves
atomic formulas, so it preserves \(\Psi\). That lemma and its proof concern
the finite sets over finitely many individuals and do not use that the
individuals form a set. Let \(\Psi\) have \(q\ge1\) prefix quantifiers,
\(q\le j\). As in the cited proof, a universal first quantifier is reduced
to an existential one with the same number of quantifiers, so let
\(\Psi=(\exists u)\Phi\), where \(\Phi\) has \(q-1\) prefix quantifiers and
free variables among \(x_1,\ldots,x_k,u\). Suppose
\(\operatorname{stsf}_{\mathfrak M_1}\Phi(x_1,\ldots,x_k,x_{k+1})\). I
chooses \(\mathfrak M_1\) and lists \(s(x_{k+1})\), or one individual if
\(s(x_{k+1})\) is empty. As \(j\ge1\), \(W_j(P)\) gives an answer \(R\)
with \(W_{j-1}(P\cup R)\); by (1), \(P\cup R\) is an isomorphism for
\(A\), and its map \(F'\) extends \(F\) by the uniqueness clause of
\cref{thm:p03-w03-p03-038-e61-hull}. Put \(y_{k+1}\defeq F'(x_{k+1})\).
Then \(P\cup R\) establishes an isomorphism between \(x_i\) and \(y_i\)
for \(i\le k+1\), and \(\Phi\) has at most \(j-1\) prefix quantifiers, so
the induction hypothesis gives
\(\operatorname{stsf}_{\mathfrak M_2}\Phi(y_1,\ldots,y_{k+1})\), and
\(\Psi\) holds of \(y_1,\ldots,y_k\). The converse direction exchanges the
models; \(W_j\) is symmetric in them. For the last claim, every closed
formula of \(\mathfrak H(A)\) is equivalent in every model to a closed
prenex formula (the prenex part of the proof of
\cref{thm:p03-w03-p03-038-e61-thm1}(3), whose equivalences are checked
clause by clause and hold in a model whose individuals form a class); if
it has \(n\) quantifiers, \(W_n(\emptyset)\) and the first claim with
\(k=0\) give agreement. The formulas of \(\mathfrak G(A)\) are formulas of
\(\mathfrak H(A)\), so the models are also elementarily indiscernible.
\end{proof}
\end{annexitem}

\begin{annexitem}{annex:p03-w03-p03-038-e61-toolkit}{thm:p03-w03-p03-038-e61-toolkit}{Ordinal arithmetic used in Part~II}
\begin{proof}
We use the recursive definitions of \cref{def:p06-ordinal-arithmetic},
the laws of \cref{thm:p06-ordinal-arithmetic-laws} (associativity, left
distributivity, \(\omega^{\beta+\gamma}=\omega^\beta\omega^\gamma\),
strict monotonicity of \(\alpha+\beta\) in \(\beta\) and of
\(\alpha\cdot\beta\) in \(\beta\) for \(\alpha>0\), weak monotonicity of
\(\omega^\beta\) in \(\beta\), continuity in the right argument), and the
existence and uniqueness of Cantor normal forms
(\cref{thm:p06-cantor-normal-form}).

\emph{(1).} Each claim is an induction on the right argument \(\xi\) or
\(\beta\). For \(\xi\le\alpha+\xi\): \(0\le\alpha\); if
\(\xi\le\alpha+\xi\), then \(\xi+1\le(\alpha+\xi)+1=\alpha+(\xi+1)\); at a
limit \(\lambda\), each \(\xi<\lambda\) satisfies
\(\xi\le\alpha+\xi\le\alpha+\lambda\), so \(\lambda\le\alpha+\lambda\).
For \(\xi\le\alpha\xi\) with \(\alpha\ge1\): \(\alpha(\xi+1)=\alpha\xi+
\alpha\ge\alpha\xi+1\ge\xi+1\), and limits as before. For
\(\alpha\le\alpha'\): \(\alpha+0\le\alpha'+0\); successor steps add \(1\)
to both sides; at limits take suprema of the inequalities. For products:
\(\alpha\cdot0=\alpha'\cdot0\); if \(\alpha\beta\le\alpha'\beta\), then
\(\alpha\beta+\alpha\le\alpha'\beta+\alpha\le\alpha'\beta+\alpha'\) by the
previous claim and right monotonicity; limits by suprema.

\emph{(2).} Uniqueness: \(\gamma<\gamma'\) gives
\(\alpha+\gamma<\alpha+\gamma'\). Existence: by (1),
\(\alpha+\beta\ge\beta\), so there is a least \(\gamma\) with
\(\alpha+\gamma\ge\beta\). If \(\gamma=0\), \(\alpha\ge\beta\), so
\(\alpha=\beta\). If \(\gamma=\gamma'+1\), \(\alpha+\gamma'<\beta\), so
\(\alpha+\gamma=(\alpha+\gamma')+1\le\beta\). If \(\gamma\) is a limit,
\(\alpha+\gamma=\sup_{\xi<\gamma}(\alpha+\xi)\le\beta\). In each case
\(\alpha+\gamma=\beta\).

\emph{(3).} By (1), \(\delta(v+1)\ge v+1>v\), so there is a least \(q'\)
with \(\delta q'>v\). It is not \(0\), and not a limit, since then
\(\delta q'=\sup_{\xi<q'}\delta\xi\le v\). So \(q'=q+1\) and
\(\delta q\le v<\delta q+\delta\). By (2), \(v=\delta q+r\) for some
\(r\), and \(r<\delta\) by strict monotonicity. If
\(\delta q+r=\delta q_2+r_2\) with \(q<q_2\), then
\(\delta q+r<\delta q+\delta=\delta(q+1)\le\delta q_2\le\delta q_2+r_2\),
which is impossible; so \(q=q_2\), and \(r=r_2\) by uniqueness in (2).

\emph{(4).} Induction on \(r\). Write \(x=\omega^{e_1}m_1+x'\). If
\(r=1\), \(x'=0\). If \(r>1\), \(x'\) is a Cantor normal form with leading
exponent \(e_2<e_1\), so by the induction hypothesis
\(x'<\omega^{e_2+1}\le\omega^{e_1}\). In both cases
\(\omega^{e_1}m_1\le x<\omega^{e_1}m_1+\omega^{e_1}=\omega^{e_1}(m_1+1)\),
and \(\omega^{e_1}(m_1+1)\le\omega^{e_1}\omega=\omega^{e_1+1}\) because
\(m_1+1<\omega\). If \(e_1<\gamma\), then \(x<\omega^{e_1+1}\le
\omega^\gamma\); if \(e_1\ge\gamma\), then
\(x\ge\omega^{e_1}m_1\ge\omega^{e_1}\ge\omega^\gamma\). The last sentence
of (4) follows, since all exponents of \(x\) are at most \(e_1\).

\emph{(5).} Let \(x,y<\omega^\gamma\), both nonzero (otherwise the sum is
the other one). Their leading exponents are below \(\gamma\) by (4); let
\(\varepsilon\) be the larger, and \(N\) a natural number exceeding both
leading coefficients. By (4), \(x,y<\omega^\varepsilon N\) (if the leading
exponent \(e\) of \(x\) is below \(\varepsilon\), then
\(x<\omega^{e+1}\le\omega^\varepsilon\)). By (1) and strict monotonicity,
\(x+y\le\omega^\varepsilon N+y<\omega^\varepsilon N+\omega^\varepsilon N=
\omega^\varepsilon(N+N)\le\omega^{\varepsilon+1}\le\omega^\gamma\). If
\(\gamma>0\), then \(1<\omega\le\omega^\gamma\), and \(\omega^\gamma=s+1\)
would give \(s,1<\omega^\gamma\) with \(s+1=\omega^\gamma\), against the
first claim; so \(\omega^\gamma\) is a limit. For absorption with
\(\gamma>0\): \(x+\omega^\gamma=\sup_{\zeta<\omega^\gamma}(x+\zeta)
\le\omega^\gamma\) by the first claim, and \(\ge\omega^\gamma\) by (1);
for \(\gamma=0\), \(x=0\). If \(c\ge1\), write \(c=1+c'\) by (2); then
\(\omega^\gamma c=\omega^\gamma+\omega^\gamma c'\) and
\(x+\omega^\gamma c=(x+\omega^\gamma)+\omega^\gamma c'=\omega^\gamma c\).

\emph{(6).} Let \(\pi\) be the sum of the common initial terms, so
\(x=\pi+\tilde x\) and \(y=\pi+\tilde y\), where \(\tilde x,\tilde y\)
are the remaining terms; by strict monotonicity \(x<y\) iff
\(\tilde x<\tilde y\). If \(\tilde x=0\ne\tilde y\), then \(x<y\). If both
are nonzero with leading terms \(\omega^em\) and \(\omega^hp\): when
\(e<h\), \(\tilde x<\omega^{e+1}\le\omega^h\le\tilde y\) by (4); when
\(e=h\) and \(m<p\), \(\tilde x<\omega^e(m+1)\le\omega^ep\le\tilde y\).
The remaining cases are the same with \(x\) and \(y\) exchanged, and if
neither list extends the other and no difference occurs, \(x=y\).

\emph{(7).} Split \(x=x_>+x_=+x_<\) into its terms of exponent \(>h\),
\(=h\) (written \(\omega^hm^*\), with \(m^*=0\) if absent) and \(<h\), and
\(y=\omega^hp_1+y'\). By (4), \(x_<<\omega^h\), so by (5)
\(x_<+\omega^hp_1=\omega^hp_1\). By associativity and distributivity,
\(x+y=x_>+\omega^h(m^*+p_1)+y'\), a Cantor normal form because the
exponents of \(y'\) are below \(h\) and those of \(x_>\) above it.

\emph{(8).} For finite \(p\ge1\), \(x\cdot p=\omega^{e_1}(m_1p)+x'\), by
induction on \(p\): \(x\cdot(p+1)=x\cdot p+x\), and (7) with leading term
\(\omega^{e_1}m_1\) of the right summand deletes \(x'\) from the left
summand and adds \(m_1\) to the coefficient \(m_1p\). Next,
\(x\cdot\omega=\sup_px\cdot p\); by (4),
\(\omega^{e_1}(m_1p)\le x\cdot p<\omega^{e_1}(m_1p+1)\), so the supremum
is \(\sup_{n<\omega}\omega^{e_1}n=\omega^{e_1}\omega=\omega^{e_1+1}\). For
\(h\ge1\) write \(h=1+h'\) by (2); then
\(x\cdot\omega^h=(x\cdot\omega)\cdot\omega^{h'}=\omega^{e_1+1}\omega^{h'}
=\omega^{e_1+h}\). Finally, by left distributivity and associativity,
\(x\cdot y=\sum_i(x\cdot\omega^{h_i})p_i\), whose summands are
\(\omega^{e_1+h_i}p_i\) for \(h_i>0\) and, when \(h_s=0\),
\(x\cdot p_s=\omega^{e_1}(m_1p_s)+x'\). The exponents
\(e_1+h_1>e_1+h_2>\cdots\) all exceed \(e_1\), which bounds the exponents
of the last block, so the sum is a Cantor normal form.
\end{proof}
\end{annexitem}

\begin{annexitem}{annex:p03-w03-p03-038-e61-sums}{thm:p03-w03-p03-038-e61-sums}{Theorem~11 and the segment lemma}
\begin{proof}
\emph{(1).} First, a position \(P\) that pairs \(L_i\) only with \(L_i'\)
and whose parts \(P_i\) are isomorphisms for \(\eta,\iota,\nu\) is itself
one: \(\eta\) holds of all individuals; two pairs in the same part are
handled by \(P_i\); for pairs \(x\leftrightarrow x'\) in part \(i\) and
\(y\leftrightarrow y'\) in part \(l\ne i\), say \(i<l\), we have
\(x\ne y\), \(x'\ne y'\), \(x<y\) and \(x'<y'\), and the reverse
relations fail on both sides. Now induct on \(j\). For \(j=0\) each
\(P_i\) is an isomorphism, hence so is \(P\). For \(j+1\), let I list a
nonempty sequence in \(L\) (the case of \(L'\) is symmetric). For each
\(i\) whose summand \(L_i\) contains some listed entries, \(W_{j+1}(P_i)\)
gives answers in \(L_i'\) to those entries with \(W_j\) of the extended
part; for the other \(i\), \(W_j(P_i)\) holds by
\cref{thm:p03-w03-p03-038-e61-wins}(1). II answers each entry by its
answer in the corresponding summand. The new position again pairs \(L_i\)
with \(L_i'\), its parts satisfy \(W_j\), and the induction hypothesis
gives \(W_j\) of it.

\emph{(2).} The individuals of \(\mathfrak M_a\) form the ordered sum
\(\mathfrak M_{a_1}\oplus\{a_1\}\oplus\mathfrak M_{a_1a_2}\oplus\cdots
\oplus\{a_k\}\oplus\mathfrak M_{a_ka}\), and those of \(\mathfrak M_b\)
the corresponding sum with the \(b_i\) (for \(b=\infty\), the last
summand is \(\mathfrak M_{b_k\infty}\)). The position
\(P=\{a_i\leftrightarrow b_i\}\) pairs \(\{a_i\}\) with \(\{b_i\}\). On
a pair of one-point summands, \(W_m(\{a_i\leftrightarrow b_i\})\) holds
for every \(m\): every move lists copies of the single element, II
answers with copies of the other, and the position stays
\(\{a_i\leftrightarrow b_i\}\), an isomorphism. On the remaining summands
the part of \(P\) is empty and \(W_{n-1}(\emptyset)\) holds by hypothesis.
Part (1) gives \(W_{n-1}(P)\), and by
\cref{thm:p03-w03-p03-038-e61-wins}(3) II wins \(H_n\) after the first
move \(P\) when both models are sets.

\emph{(3).} The individuals of \(\mathfrak M_{a+d}\) form
\(\mathfrak M_a\oplus\{x:a\le x<a+d\}\), and the second summand is
isomorphic to the ordinal \(d\) by \(a+t\mapsto t\)
(\cref{thm:p03-w03-p03-038-e61-toolkit}(2) and strict monotonicity);
likewise for \(b+d\). Between two copies of \(d\), \(W_n(\emptyset)\)
holds because II can answer each entry by the same element, by induction
on the number of moves; transported by
\cref{thm:p03-w03-p03-038-e61-wins}(2), \(W_n(\emptyset)\) holds between
the second summands. Part (1) with \(P=\emptyset\) gives
\(W_n(\mathfrak M_{a+d},\mathfrak M_{b+d})\).
\end{proof}
\end{annexitem}

\begin{annexitem}{annex:p03-w03-p03-038-e61-thm12}{thm:p03-w03-p03-038-e61-thm12}{Theorem~12 of the source and its class form}
\begin{proof}
All references to (1)--(8) are to
\cref{thm:p03-w03-p03-038-e61-toolkit}. By
\cref{thm:p03-w03-p03-038-e61-wins}(2), \(W_j\) for a pair of segments
depends only on their order types; by
\cref{def:p03-w03-p03-038-e61-congruence}, an open segment of type \(t\)
is isomorphic to \(\mathfrak M_t\), and \(\mathfrak M_{x\infty}\) to
\(\mathfrak M_\infty\).

\emph{Construction.} Let \(k\ge1\) and \(x_1<\cdots<x_l\) be ordinals.
Let \(v_1\defeq x_1\) and \(v_i\) the type of
\(\{z:x_{i-1}<z<x_i\}\) for \(2\le i\le l\). By (3),
\(v_i=\omega^kq_i+r_i\) with \(r_i<\omega^k\). Define
\(y_1\defeq\omega^k\operatorname{sgn}q_1+r_1\) and
\(y_i\defeq y_{i-1}+1+\omega^k\operatorname{sgn}q_i+r_i\). Then:
(a) \(y_1<\cdots<y_l\) by strict monotonicity, and \(y_l<\omega^{k+1}\)
by (5), since every summand \(1\) or \(\omega^k\operatorname{sgn}q_i+r_i
<\omega^k+\omega^k\le\omega^{k+1}\) is below \(\omega^{k+1}\); (b) the
initial segment below \(y_1\) has type \(y_1\), and the open segment
between \(y_{i-1}\) and \(y_i\) has type
\(\omega^k\operatorname{sgn}q_i+r_i\) by uniqueness in (2); this is
\(\equiv v_i\pmod{\omega^k}\), being equal to \(v_i\) when \(q_i=0\), and
of the form \(\omega^k\cdot1+r_i\) against \(\omega^kq_i+r_i\) with
\(q_i\ne0\) otherwise; (c) if \(e=\omega^{k+1}c'\) with \(c'\ge1\), then
\(y_l+1<\omega^{k+1}\) by (a) and (5), so \((y_l+1)+e=e\) by (5), and the
open segment between \(y_l\) and \(e\) has type
\(e=\omega^k\cdot(\omega c')\).

\emph{The last segment.} Let \(a=\omega^{k+1}c\) with \(c\ge1\) and
\(x<a\). By (3), the remainder of \(a=\omega\cdot(\omega^kc)\) on division
by \(\omega\) is \(0\), while a successor \(s+1\), with
\(s=\omega q+r\), equals \(\omega q+(r+1)\) with remainder \(r+1\ne0\);
so \(a\) is a nonzero limit and \(x+1<a\). Let \(t\) be the type of the
open segment between \(x\) and \(a\), so \(t>0\), and write
\(t=\omega^kq+r\), \(r<\omega^k\). If \(r\ne0\), divide
\(s=(x+1)+\omega^kq\) as \(\omega^kq''+r''\); then
\(a=\omega^kq''+(r''+r)\) with \(0<r''+r<\omega^k\) by (5), against the
remainder \(0\) of \(a=\omega^k(\omega c)\). So \(t=\omega^kq\) with
\(q\ge1\).

\emph{Induction on \(n\).} We prove (1) and (2) together. For \(n=1\):
if \(a=b\), II answers each entry by itself. Otherwise in (1),
\(a=\omega d+e\ge\omega\) and \(b\ge\omega\), and in (2), \(a\ge\omega\)
and \(\mathfrak M_\infty\) is infinite. Given a nonempty list in one
model, let \(c_0<\cdots<c_{p-1}\) be its distinct entries; II answers
each entry \(c_i\) by the natural number \(i\), an individual of the other
model. The resulting position respects \(\iota\) and \(\nu\), so
\(W_1(\emptyset)\).

Assume (1) and (2) for \(n=k\ge1\). \emph{(1) for \(k+1\).} If \(a=b\),
II copies. Otherwise \(a=\omega^{k+1}c+d\), \(b=\omega^{k+1}c'+d\) with
\(c,c'\ne0\); by \cref{thm:p03-w03-p03-038-e61-sums}(3) it suffices to
treat \(d=0\). The hypothesis is symmetric in \(a\) and \(b\), so let I
list entries in \(\mathfrak M_a\), with distinct values
\(a_1<\cdots<a_l\). Apply the construction to them, obtaining
\(b_1<\cdots<b_l<\omega^{k+1}\le b\), and let II answer each entry
\(a_i\) by \(b_i\). By (b) and the induction hypothesis (1), \(W_k\) holds
for the initial segments and for each pair of open segments between
consecutive points. For the last pair, the segment of \(\mathfrak M_a\)
above \(a_l\) has type \(\omega^kq\) with \(q\ge1\), and that of
\(\mathfrak M_b\) above \(b_l\) has type \(\omega^k(\omega c')\) by (c);
these are congruent modulo \(\omega^k\) with \(e=0\), so \(W_k\) holds
for them. \Cref{thm:p03-w03-p03-038-e61-sums}(2) gives \(W_k\) of the
position \(\{a_i\leftrightarrow b_i\}\). Every first move has such an
answer, so \(W_{k+1}(\mathfrak M_a,\mathfrak M_b)\).

\emph{(2) for \(k+1\).} Let \(a=\omega^{k+1}c\), \(c\ge1\). If I lists
entries in \(\mathfrak M_a\) with distinct values \(a_1<\cdots<a_l\), II
answers by the \(b_i\) of the construction. The segments below and
between the points are handled by (b) and the induction hypothesis (1).
The last pair is the segment of \(\mathfrak M_a\) above \(a_l\), of type
\(\omega^kq\) with \(q\ge1\), against \(\mathfrak M_{b_l\infty}\cong
\mathfrak M_\infty\); \(W_k\) holds for it by the induction hypothesis (2).
If I lists entries in \(\mathfrak M_\infty\), with distinct values
\(b_1<\cdots<b_l\), apply the construction to the \(b_i\), obtaining
\(a_1<\cdots<a_l<\omega^{k+1}\le a\), and let II answer \(b_i\) by
\(a_i\). The segments below and between the points are handled as
before. The segment of \(\mathfrak M_a\) above \(a_l\) has type
\(a=\omega^k(\omega c)\) by (c), and the segment of
\(\mathfrak M_\infty\) above \(b_l\) is isomorphic to
\(\mathfrak M_\infty\); \(W_k\) holds for them by the induction hypothesis
(2). \Cref{thm:p03-w03-p03-038-e61-sums}(2), with \(b=\infty\), gives
\(W_k\) of the position, hence \(W_{k+1}(\mathfrak M_a,
\mathfrak M_\infty)\).

For set models, \cref{thm:p03-w03-p03-038-e61-wins}(3) turns
\(W_n(\mathfrak M_a,\mathfrak M_b)\) into a winning strategy for II in
\(H_n(\mathfrak M_a,\mathfrak M_b)\).
\end{proof}
\end{annexitem}

\begin{annexitem}{annex:p03-w03-p03-038-e61-thm13}{thm:p03-w03-p03-038-e61-thm13}{Theorem~13 of the source}
\begin{proof}
Since \(n+\omega=\sup_{m<\omega}(n+m)=\omega\), the exponent law gives
\(\omega^n\omega^\omega=\omega^\omega\) for every natural number \(n\).
If \(a=b\), the models coincide. Otherwise \(a=\omega^\omega d+e\) and
\(b=\omega^\omega d'+e\) with \(d,d'\ne0\), so
\(a=\omega^n(\omega^\omega d)+e\) and \(b=\omega^n(\omega^\omega d')+e\)
by associativity, with nonzero coefficients; hence
\(a\equiv b\pmod{\omega^n}\) for every \(n\ge1\), and
\cref{thm:p03-w03-p03-038-e61-thm12}(1) gives
\(W_n(\mathfrak M_a,\mathfrak M_b)\) for every \(n\ge1\), while
\(W_0(\emptyset)\) holds trivially. If \(a=0\), then \(b=0\); otherwise
both are models, and \cref{thm:p03-w03-p03-038-e61-wins}(4) gives
indiscernibility by means of finite sets and elementary
indiscernibility. For the class form, \(\omega^\omega c=
\omega^n(\omega^\omega c)\) with \(\omega^\omega c\ge1\), so
\cref{thm:p03-w03-p03-038-e61-thm12}(2) gives
\(W_n(\mathfrak M_{\omega^\omega c},\mathfrak M_\infty)\) for every
\(n\ge1\), and \cref{thm:p03-w03-p03-038-e61-wins}(4) applies.
\end{proof}
\end{annexitem}

\begin{annexitem}{annex:p03-w03-p03-038-e61-lift}{thm:p03-w03-p03-038-e61-lift}{Lifting a correspondence of exponents}
\begin{proof}
References (1)--(8) are to \cref{thm:p03-w03-p03-038-e61-toolkit}.

\emph{(1).} As \(f\) preserves \(<\), the exponents
\(f(e_1)>f(e_2)>\cdots\) are decreasing, so \(g_f(x)\) is written in Cantor
normal form, with exponents in the range of \(f\); and
\(g_{f^{-1}}(g_f(x))=x\), \(g_f(g_{f^{-1}}(y))=y\). Order: by (6),
whether \(x<y\) depends only on the first place where the term lists
differ and on comparisons of exponents and coefficients there; \(f\)
preserves equality and order of exponents, so the term lists of
\(g_f(x)\) and \(g_f(y)\) differ first at the same place, in the same way.
Hence \(x<y\) implies \(g_f(x)<g_f(y)\); the converse is the same claim
for \(f^{-1}\), applied to \(g_f(x),g_f(y)\). Addition: if \(y=0\), both
sides are trivial. Otherwise, by (7), the Cantor normal form of \(x+y\) is
determined by the terms of \(x\) and \(y\) and by comparisons of the
exponents of \(x\) with the leading exponent of \(y\); its exponents are
exponents of \(x\) or \(y\), so \(x+y\in E(f)\), and the same recipe
applied to \(g_f(x)\) and \(g_f(y)\) yields \(g_f(x+y)\). So
\(g_f(x)+g_f(y)=g_f(x+y)\). If \(x+y=z\), then
\(g_f(x)+g_f(y)=g_f(z)\); if \(g_f(x)+g_f(y)=g_f(z)\), then
\(g_f(x+y)=g_f(z)\) and \(x+y=z\), because \(g_f\) is injective.

\emph{(2).} The hypothesis also holds for \(f^{-1}\), so it suffices to
show that \(x\cdot y=z\) implies \(g_f(x)\cdot g_f(y)=g_f(z)\) for
\(x,y,z\in E(f)\); the converse is this claim for \(f^{-1}\) applied to
\(g_f(x),g_f(y),g_f(z)\in E(f^{-1})\). If \(x=0\) or \(y=0\), then
\(z=0\) and both sides are \(0\). Otherwise let \(x\) have leading term
\(\omega^{e_1}m_1\) and remaining terms \(x'\), and
\(y=\sum_{i\le s}\omega^{h_i}p_i\). For \(u\) in the domain of \(f\):
\(f(u)=0\) iff \(u=0\), since \(u+u=u\) holds only for \(u=0\) and \(f\)
preserves and reflects this equation. By (8), the Cantor normal form of
\(z=x\cdot y\) has the terms \(\omega^{e_1+h_i}p_i\) for \(h_i>0\),
followed, if \(h_s=0\), by \(\omega^{e_1}(m_1p_s)+x'\). As \(z\in E(f)\),
each \(e_1+h_i\) lies in the domain of \(f\), and
\(f(e_1)+f(h_i)=f(e_1+h_i)\). By (8) again, \(g_f(x)\cdot g_f(y)\) has the
terms \(\omega^{f(e_1)+f(h_i)}p_i=\omega^{f(e_1+h_i)}p_i\) for
\(f(h_i)>0\), that is for \(h_i>0\), followed, if \(f(h_s)=0\), by
\(\omega^{f(e_1)}(m_1p_s)+g_f(x')\). This is \(g_f(z)\).

\emph{(3).} Each sum \(\omega^{a_1}n_1+\cdots+\omega^{a_k}n_k\) with
\(a_i<a\) is below \(\omega^a\): each term is below
\(\omega^{a_i+1}\le\omega^a\) by (4), and finite sums stay below
\(\omega^a\) by (5). Every \(x<\omega^a\) is such a sum, namely its Cantor
normal form, whose exponents are below \(a\) by (4); the empty sum gives
\(0\). So these sums are exactly the ordinals below \(\omega^a\), and the
least ordinal greater than all of them is \(\omega^a\).
\end{proof}
\end{annexitem}

\begin{annexitem}{annex:p03-w03-p03-038-e61-thm14}{thm:p03-w03-p03-038-e61-thm14}{Theorems~14 and~16 of the source}
\begin{proof}
\emph{(1).} For a position \(Q\) of the pair
\((\mathfrak M_a,\mathfrak M_b)\) that is an isomorphism for \(A_1\), let
\(f_Q\) be the map it defines; it is a finite injective map preserving
\(<\). We show by induction on \(j\): if \(W_j(Q)\) for
\((\mathfrak M_a,\mathfrak M_b)\) over \(\mathcal T_1\), and \(P\) is a
position of \((\mathfrak M_{\omega^a},\mathfrak M_{\omega^b})\) each of
whose pairs has the form \(x\leftrightarrow g_{f_Q}(x)\) with
\(x\in E(f_Q)\), then \(W_j(P)\) over \(\mathcal T_2\). By
\cref{thm:p03-w03-p03-038-e61-wins}(1), \(Q\) is an isomorphism, so
\(f_Q\) is defined. For \(j=0\), \(P\) is an isomorphism for \(\eta\),
\(\iota\), \(\nu\) and \(\sigma\) by
\cref{thm:p03-w03-p03-038-e61-lift}(1). For \(j+1\), let I list
\(z_1,\ldots,z_p\) in \(\mathfrak M_{\omega^a}\), and let \(X\) be the set
of exponents of their Cantor normal forms, all below \(a\) by
\cref{thm:p03-w03-p03-038-e61-toolkit}(4). If \(X\) is empty, every
\(z_i\) is \(0=g_{f_Q}(0)\); II answers \(0\), and \(W_j(Q)\) holds by
\cref{thm:p03-w03-p03-038-e61-wins}(1). Otherwise I's listing of \(X\)
in \(\mathfrak M_a\) is a move in the game on
\((\mathfrak M_a,\mathfrak M_b)\); \(W_{j+1}(Q)\) gives an answer \(R\)
with \(W_j(Q\cup R)\). Let \(f'=f_{Q\cup R}\), which extends \(f_Q\);
then \(g_{f'}\) extends \(g_{f_Q}\), and II answers each \(z_i\) by
\(g_{f'}(z_i)\), which lies below \(\omega^b\) by
\cref{thm:p03-w03-p03-038-e61-toolkit}(4) as its exponents lie below
\(b\). If I lists \(z_1,\ldots,z_p\) in \(\mathfrak M_{\omega^b}\), the
same is done with the exponents listed in \(\mathfrak M_b\): after the
answer \(R\), the exponents of each \(z_i\) lie in the range of \(f'\),
and II answers \(z_i\) by \(g_{f'^{-1}}(z_i)\), whose image under
\(g_{f'}\) is \(z_i\) by \cref{thm:p03-w03-p03-038-e61-lift}(1), and
which lies below \(\omega^a\) because its exponents lie in the domain of
\(f'\), below \(a\). In each
case the new position consists of pairs \(x\leftrightarrow g_{f'}(x)\),
and the induction hypothesis gives \(W_j\) of it. Taking
\(Q=P=\emptyset\) and \(j=n\) proves (1).

\emph{(2).} The same induction with \(\mathcal T_2\) and \(\mathcal T_3\)
in place of \(\mathcal T_1\) and \(\mathcal T_2\). Now \(Q\) is an
isomorphism for \(\nu\) and \(\sigma\), so \(f_Q\) preserves \(<\), and
\(u+v=w\) iff \(f_Q(u)+f_Q(v)=f_Q(w)\) on its domain; by
\cref{thm:p03-w03-p03-038-e61-lift}(1) and (2), a position of pairs
\(x\leftrightarrow g_{f_Q}(x)\) is an isomorphism for \(\nu\), \(\sigma\)
and \(\mu\).

\emph{(3).} Replace \(\mathfrak M_b\) and \(\mathfrak M_{\omega^b}\) by
\(\mathfrak M_\infty\). Exponents of any ordinal are ordinals, so the
moves of the inner game are legal, and no bound on the answers is needed.
The induction is otherwise unchanged.
\end{proof}
\end{annexitem}

\begin{annexitem}{annex:p03-w03-p03-038-e61-thm15}{thm:p03-w03-p03-038-e61-thm15}{Theorems~15 and~17 of the source}
\begin{proof}
As in the proof of \cref{thm:p03-w03-p03-038-e61-thm13},
\(a\equiv b\pmod{\omega^\omega}\) gives
\(W_n(\mathfrak M_a,\mathfrak M_b)\) over \(\mathcal T_1\) for every
\(n\). \emph{(1).} By \cref{thm:p03-w03-p03-038-e61-thm14}(1),
\(W_n(\mathfrak M_{\omega^a},\mathfrak M_{\omega^b})\) over
\(\mathcal T_2\) for every \(n\); these are models, since
\(\omega^a,\omega^b\ge1\), and
\cref{thm:p03-w03-p03-038-e61-wins}(4) gives the claim. \emph{(2).}
Applying \cref{thm:p03-w03-p03-038-e61-thm14}(2) to the pair of (1) gives
\(W_n(\mathfrak M_{\omega^{\omega^a}},\mathfrak M_{\omega^{\omega^b}})\)
over \(\mathcal T_3\) for every \(n\), and the same lemma applies.
\emph{(3).} By \cref{thm:p03-w03-p03-038-e61-thm12}(2) with
\(\omega^\omega=\omega^n\omega^\omega\),
\(W_n(\mathfrak M_{\omega^\omega},\mathfrak M_\infty)\) over
\(\mathcal T_1\) for every \(n\ge1\). By
\cref{thm:p03-w03-p03-038-e61-thm14}(3), first with (1) and then with
(2), \(W_n(\mathfrak M_{\omega^{\omega^\omega}},\mathfrak M_\infty)\) over
\(\mathcal T_2\) and
\(W_n(\mathfrak M_{\omega^{\omega^{\omega^\omega}}},\mathfrak M_\infty)\)
over \(\mathcal T_3\), for every \(n\);
\cref{thm:p03-w03-p03-038-e61-wins}(4) concludes.
\end{proof}
\end{annexitem}

\begin{annexitem}{annex:p03-w03-p03-038-e61-converse}{thm:p03-w03-p03-038-e61-converse}{The ordinals \(\omega\) and \(\omega+1\)}
\begin{proof}
Dividing by \(\omega\), \(\omega=\omega\cdot1+0\) and
\(\omega+1=\omega\cdot1+1\); by uniqueness in
\cref{thm:p03-w03-p03-038-e61-toolkit}(3), no ordinal \(e\) is the
remainder of both, and \(\omega\ne\omega+1\), so
\(\omega\not\equiv\omega+1\pmod\omega\). \(W_1\) holds because both models
are infinite: II answers the \(i\)th smallest distinct entry of a list by
the natural number \(i\), as in the base case of the proof of
\cref{thm:p03-w03-p03-038-e61-thm12}. \(W_2\) fails: let I list
\(\omega\) in \(\mathfrak M_{\omega+1}\). Any answer is a natural number
\(m\). Then I lists \(m+1\) in \(\mathfrak M_\omega\). An answer \(y\) of
II satisfies \(y\le\omega\); if \(y=\omega\), the pairs
\(m\leftrightarrow\omega\), \(m+1\leftrightarrow\omega\) violate
\(\iota\); if \(y<\omega\), then \(m<m+1\) while \(\omega<y\) fails, which
violates \(\nu\). So no answer to the second move keeps \(W_0\), hence
\(W_1(\{m\leftrightarrow\omega\})\) fails for every \(m\), and
\(W_2(\emptyset)\) fails.
\end{proof}
\end{annexitem}

\begin{annexitem}{annex:p03-w03-p03-038-e61-below}{thm:p03-w03-p03-038-e61-below}{Ordinals below \(a_i\) are definable}
\begin{proof}
All formulas below lie in \(\mathfrak G(A_i)\): their quantifiers are
restricted to \(\eta\), written \((\exists w)_\eta\) and
\((\forall w)_\eta\). Each defining formula begins with the conjunct
\(\eta(v)\), so it fails at every element of \(\mathfrak M_\infty\) that is
not an individual. References (1)--(8) are to
\cref{thm:p03-w03-p03-038-e61-toolkit}.

\emph{Multiples of \(\omega^k\) in \(T_1\).} Let
\(Z(v)\defeq(\forall w)_\eta\sim\nu(w,v)\), true of \(0\) only. Let
\(D_0(v)\defeq\iota(v,v)\) and
\[D_{k+1}(v)\defeq Z(v)\vee(\forall w)_\eta\bigl(\nu(w,v)\supset
(\exists z)_\eta(\nu(w,z)\wedge\nu(z,v)\wedge D_k(z))\bigr).\]
We show that \(D_k(x)\) holds in \(\mathfrak M_\infty\) iff
\(x=\omega^kq\) for some \(q\). For \(k=0\) every \(x=1\cdot x\). Assume it
for \(k\), and let \(x\ne0\); then \(D_{k+1}(x)\) says that multiples of
\(\omega^k\) occur strictly between every \(y<x\) and \(x\).
If \(x=\omega^{k+1}q'=\omega^k(\omega q')\) with \(q'\ge1\) and \(y<x\),
write \(y=\omega^ks+r\), \(r<\omega^k\), by (3). Then \(s<\omega q'\),
since otherwise \(y\ge\omega^ks\ge x\) by (1); \(\omega q'\) is a limit,
since its remainder on division by \(\omega\) is \(0\) while every
successor has a nonzero remainder (the argument for \(a\) in the proof of
\cref{thm:p03-w03-p03-038-e61-thm12}); so \(s+1<\omega q'\), and
\(z=\omega^k(s+1)=\omega^ks+\omega^k\) satisfies \(y<z<x\).
Conversely, let multiples of \(\omega^k\) occur arbitrarily high below
\(x\ne0\). Write \(x=\omega^{k+1}q'+r'\) with \(r'<\omega^{k+1}\), and
\(r'=\omega^kj+r''\) with \(r''<\omega^k\); here \(j<\omega\), since
\(j\ge\omega\) would give \(r'\ge\omega^k\omega\). If \(r''\ne0\), let
\(y=\omega^{k+1}q'+\omega^kj<x\). Every \(z\) with \(y<z<x\) is
\(y+t\) with \(0<t<r''<\omega^k\) by (2), so
\(z=\omega^k(\omega q'+j)+t\) has nonzero remainder on division by
\(\omega^k\), and by uniqueness in (3) it is not a multiple of
\(\omega^k\): a contradiction. If \(r''=0\) and \(j\ge1\), let
\(y=\omega^k(\omega q'+(j-1))\); then \(x=y+\omega^k\), and every \(z\)
strictly between is \(y+t\) with \(0<t<\omega^k\), again not a multiple.
So \(r'=0\), and \(x=\omega^{k+1}q'\) with \(q'\ge1\) as \(x\ne0\).

\emph{The next multiple.} Let
\(N_k(u,v)\defeq D_k(v)\wedge\nu(u,v)\wedge(\forall w)_\eta
((D_k(w)\wedge\nu(u,w))\supset\sim\nu(w,v))\). If \(u=\omega^ks\), the
only \(v\) with \(N_k(u,v)\) is \(\omega^k(s+1)=u+\omega^k\): it is a
multiple above \(u\), and a multiple \(\omega^kt>u\) has \(t>s\), so
\(t\ge s+1\) and \(\omega^kt\ge\omega^k(s+1)\).

\emph{\(i=1\).} \(0\) is defined by \(\varphi_0(v)\defeq\eta(v)\wedge
Z(v)\). If \(y\) is defined by \(\varphi_y\) and is a multiple of
\(\omega^k\), then \(y+\omega^k\) is defined by
\(\eta(v)\wedge(\exists u)_\eta(\varphi_y(u)\wedge N_k(u,v))\). Let
\(0<x<\omega^\omega\), with Cantor normal form
\(\omega^{k_1}m_1+\cdots+\omega^{k_r}m_r\), \(k_1>\cdots>k_r\) natural
numbers. Starting from \(0\), add \(\omega^{k_1}\) \(m_1\) times, then
\(\omega^{k_2}\) \(m_2\) times, and so on. Before each addition of
\(\omega^{k_l}\), the current value is a sum of terms \(\omega^{k_i}\)
with \(k_i\ge k_l\), and \(\omega^{k_i}=\omega^{k_l}\omega^{k_i-k_l}\), so
by distributivity it is a multiple of \(\omega^{k_l}\). Each step
therefore preserves definability, and \(x\) is definable in \(T_1\).

\emph{Additive principal ordinals.} Let
\[
\begin{aligned}
\mathrm{AP}(v)\defeq{}&\eta(v)\wedge\sim Z(v)\
&\wedge(\forall y)_\eta(\forall z)_\eta(\forall w)_\eta
((\nu(y,v)\wedge\nu(z,v)\wedge\sigma(y,z,w))\supset\nu(w,v)).
\end{aligned}
\]
It holds
of \(x\) iff \(x=\omega^\gamma\) for some \(\gamma\). If
\(x=\omega^\gamma\), use (5). If \(x\ne0\) is not a power of \(\omega\),
write \(x=\omega^em+x'\) in Cantor normal form with \(m\ge2\) or
\(x'\ne0\). If \(m\ge2\), take \(y=\omega^e\) and
\(z=\omega^e(m-1)+x'\); both are below \(x\) (by strict monotonicity and
\(x'<\omega^e\), from (4)), and \(y+z=x\). If \(m=1\) and \(x'\ne0\), take
\(y=\omega^e\), \(z=x'\); again \(y,z<x\) and \(y+z=x\). The map
\(\gamma\mapsto\omega^\gamma\) is strictly increasing (by (4),
\(\gamma<\gamma'\) gives \(\omega^\gamma<\omega^{\gamma+1}\le
\omega^{\gamma'}\)) and onto the ordinals satisfying \(\mathrm{AP}\); it
carries \(+\) to \(\cdot\), because
\(\omega^\gamma\omega^{\gamma'}=\omega^{\gamma+\gamma'}\).

\emph{Relativization.} For \(\varphi\in T_1\) (resp.\ \(T_2\)) let
\(\varphi^*\) replace each \((\exists u)_\eta\psi\) by
\((\exists u)_\eta(\mathrm{AP}(u)\wedge\psi)\) and each
\((\forall u)_\eta\psi\) by \((\forall u)_\eta(\mathrm{AP}(u)\supset\psi)\),
and, for \(T_2\), each atomic \(\sigma(\cdot)\) by \(\mu(\cdot)\). By
induction on \(\varphi\): for ordinals \(\gamma_1,\ldots,\gamma_m\),
\(\varphi(\gamma_1,\ldots,\gamma_m)\) holds in \(\mathfrak M_\infty\) iff
\(\varphi^*(\omega^{\gamma_1},\ldots,\omega^{\gamma_m})\) does. Atomic
\(\eta\), \(\iota\) and \(\nu\) are preserved and reflected by an
injective increasing map; \(\sigma(\gamma,\gamma',\gamma'')\) holds iff
\(\mu(\omega^\gamma,\omega^{\gamma'},\omega^{\gamma''})\) does by the
exponent law and injectivity; connectives commute with the translation;
and a restricted quantifier over all ordinals corresponds to one over the
ordinals satisfying \(\mathrm{AP}\), which are exactly the
\(\omega^\gamma\). \(\varphi^*\) lies in \(T_2\) (resp.\ \(T_3\)), as
\(\mathrm{AP}\) uses \(\sigma\).

\emph{\(i=2,3\).} If \(e\) is defined in \(T_1\) (resp.\ \(T_2\)) by
\(\varphi_e\), then \(\omega^e\) is defined in \(T_2\) (resp.\ \(T_3\)) by
\(\mathrm{AP}(v)\wedge\varphi_e^*(v)\): an element satisfies it iff it is
an individual \(\omega^\gamma\) with \(\varphi_e(\gamma)\), that is
\(\gamma=e\). If \(p\) and \(q\) are defined by \(\Psi_1\) and \(\Psi_2\),
then \(p+q\) is defined by \(\eta(v)\wedge(\exists u)_\eta(\exists w)_\eta
(\Psi_1(u)\wedge\Psi_2(w)\wedge\sigma(u,w,v))\). Let \(0<x<a_2\). By (4),
the exponents of its Cantor normal form are below \(\omega^\omega=a_1\),
so each is definable in \(T_1\) by the case \(i=1\); hence each
\(\omega^{e_j}\) is definable in \(T_2\), and so is \(x\), a finite sum of
them. Likewise, for \(0<x<a_3\) the exponents are below \(a_2\), hence
definable in \(T_2\), so each \(\omega^{e_j}\), and \(x\), is definable in
\(T_3\). The case \(x=0\) is \(\varphi_0\), which lies in every \(T_i\).

\emph{(\(\bullet\bullet\)).} A formula of \(T_i=\mathfrak G(A_i)\) is a
formula of \(\mathcal T_i=\mathfrak H(A_i)\), with the same satisfaction
in every model, so a defining formula in \(T_i\) is one in
\(\mathcal T_i\).
\end{proof}
\end{annexitem}

\begin{annexitem}{annex:p03-w03-p03-038-e61-initial}{thm:p03-w03-p03-038-e61-initial}{The initial-segment lemma}
\begin{proof}
Write \(\mathfrak M=\mathfrak M_\infty\) and \(D=D(\mathfrak M,
\mathcal T_i)\). For \(x\in D\) and a formula \(\Psi\) defining \(x\) in
\(\mathfrak M\), the closed formulas
\((\exists v)(\eta(v)\wedge\Psi(v))\) and
\((\forall v)(\forall w)((\Psi(v)\wedge\Psi(w))\supset\iota(v,w))\) hold
in \(\mathfrak M\), hence in \(\mathfrak M_a\); so exactly one element of
\(\mathfrak M_a\) satisfies \(\Psi\), and it is an individual. Call it
\(h(x)\). It does not depend on \(\Psi\): if \(\Psi'\) also defines
\(x\), then \((\forall v)(\Psi(v)\equiv\Psi'(v))\) holds in
\(\mathfrak M\), hence in \(\mathfrak M_a\). If \(x<y\) in \(D\), defined
by \(\Psi\) and \(\Psi'\), then
\((\exists u)(\exists v)(\Psi(u)\wedge\Psi'(v)\wedge\nu(u,v))\) holds in
\(\mathfrak M\), hence in \(\mathfrak M_a\), so \(h(x)<h(y)\). Thus
\(h\) maps \(D\) strictly increasingly into \(\{z:z<a\}\).

By hypothesis \(\{z:z<a\}\subseteq D\). For \(z<a\), \(h(z)\ge z\):
otherwise take the least \(z<a\) with \(h(z)<z\); then \(h(z)<a\) and
\(h(h(z))<h(z)\), against the minimality of \(z\). Suppose some
\(x\in D\) has \(x\ge a\). Then \(h(x)>h(z)\ge z\) for every \(z<a\), so
\(h(x)\ge a\), contradicting \(h(x)<a\). Hence \(D=\{z:z<a\}\).
\end{proof}
\end{annexitem}

\begin{annexitem}{annex:p03-w03-p03-038-e61-thm18}{thm:p03-w03-p03-038-e61-thm18}{Theorem~18 of the source}
\begin{proof}
(c) implies (b) by \cref{thm:p03-w03-p03-038-e61-below}(\(\bullet\)), and
(b) implies (a) by (\(\bullet\bullet\)). For (a) implies (c): by
(\(\bullet\)) and (\(\bullet\bullet\)), every \(x<a_i\) lies in
\(D(\mathfrak M_\infty,\mathcal T_i)\). \(\mathfrak M_{a_i}\) and
\(\mathfrak M_\infty\) are indiscernible by means of finite sets as models
for \(\mathcal T_i\): for \(i=1\) by \cref{thm:p03-w03-p03-038-e61-thm13}
with \(c=1\), and for \(i=2,3\) by
\cref{thm:p03-w03-p03-038-e61-thm15}(3). \Cref{thm:p03-w03-p03-038-e61-initial}
with \(a=a_i\) gives \(D(\mathfrak M_\infty,\mathcal T_i)=\{x:x<a_i\}\),
so every \(x\) definable in \(\mathcal T_i\) is below \(a_i\).
\end{proof}
\end{annexitem}

\begin{annexitem}{annex:p03-w03-p03-038-closed}{thm:p03-w03-p03-038-closed}{The winning condition is closed}
\begin{proof}
\emph{(1).} Let \(\xi_0\) be the least stage whose position fails to be a
partial isomorphism; it exists because the ordinals are well ordered, and
\(\xi_0\ne0\) by hypothesis. If \(\xi_0\) were a limit, the failure would be
witnessed by one atomic formula at finitely many coordinates. Each of these
belongs to the initial position or to a round before \(\xi_0\); let \(\eta\)
be the largest of these finitely many rounds. Then \(\eta+1<\xi_0\) and the
position at stage \(\eta+1\) already contains the failing coordinates,
contradicting the minimality of \(\xi_0\).

\emph{(2).} If a play is won by Spoiler, some finite stage \(n\) fails. Every
play agreeing with it in its first \(n\) rounds has the same position at
stage \(n\) and is also won by Spoiler. These plays form a basic open set
containing the given play, so Spoiler's set is open and Duplicator's is
closed.

\emph{(3).} Say that Spoiler wins from a position \(P\) if he has a strategy
such that every play from \(P\) following it is lost by Duplicator. Let
\(L\) be the set of positions from which Spoiler does not win. Let
\(P\in L\). Then \(P\) is a partial isomorphism, since otherwise every play
from \(P\) is lost by Duplicator and any strategy wins for Spoiler. For each
move \(c\) of Spoiler there is an answer \(d\) with \(P(c,d)\in L\): otherwise,
choosing by the Axiom of Choice a winning Spoiler strategy \(\rho_d\) from
\(P(c,d)\) for every answer \(d\), the strategy ``play \(c\), then follow
\(\rho_d\)'' would win from \(P\). Fix such answers by a choice function.
Duplicator's strategy from \(P\in L\) answers with them; every position of
the play then lies in \(L\), so every finite stage is a partial
isomorphism, and Duplicator wins. Hence from every position one of the
players has a winning strategy.
\end{proof}
\end{annexitem}

\begin{annexitem}{annex:p03-w03-p03-038-omega-comparison}{thm:p03-w03-p03-038-omega-comparison}{The game of length \(\omega\) and \(L_{\infty\omega}\)}
\begin{proof}
If Duplicator wins from \((\bar a,\bar b)\), let \(I\) be the set of
positions from which she has a winning strategy. It contains
\((\bar a,\bar b)\). Each member is a partial isomorphism, being stage
\(0\) of the plays from it. If \(\sigma\) wins from \(P\in I\) and Spoiler
moves \(c\) (or \(d\)), \(\sigma\)'s answer leads to a position from which
the continuation of \(\sigma\) wins, so forth and back hold. Conversely, if
a back-and-forth family \(I\) contains \((\bar a,\bar b)\), fix by the Axiom
of Choice, for every \(P\in I\) and move, an answer leading into \(I\).
Following these answers keeps every finite stage in \(I\); each is a
partial isomorphism, so Duplicator wins.

For the empty initial position, read each position of a back-and-forth
family as the finite map \(a_i\mapsto b_i\), which is well defined and
injective because the position respects equality. The maps form a nonempty
family of finite partial isomorphisms with the forth and back conditions of
\cref{def:p08-bvb501-potential-isomorphism}. Conversely, from a potential
isomorphism, closed under restrictions as permitted there, the positions
listing its maps form a back-and-forth family: if \(f\) extends to \(g\) in
the family with \(c\in\dom g\), then
\(f\cup\{(c,g(c))\}\) is a restriction of \(g\). The last equivalence is
\cref{thm:p08-bvb501-potential-isomorphism}.
\end{proof}
\end{annexitem}

\begin{annexitem}{annex:p03-w03-p03-038-clock-comparison}{thm:p03-w03-p03-038-clock-comparison}{Clock games and ordinal approximants}
\begin{proof}
Write \(W^\alpha(P)\) when Duplicator has a winning strategy in \(G^\alpha\)
from the position \(P\). \emph{Unfolding:} for \(\alpha>0\), \(W^\alpha(P)\)
holds iff \(P\) is a partial isomorphism and, for every \(\beta<\alpha\), every
move \(c\) in \(\mathcal A\) has an answer \(d\) with \(W^\beta(P(c,d))\) and
every move \(d\) in \(\mathcal B\) has an answer \(c\) with \(W^\beta(P(c,d))\).
From left to right, \(P\) is a restriction of the final position of any play
won by the strategy, and the answer of a winning strategy to the move
\((\beta,c)\) leads to a position won by its continuation. From right to left,
fix answers and continuation strategies by the Axiom of Choice. For
\(\alpha=0\), \(W^0(P)\) iff \(P\) is a partial isomorphism.

We show \(W^\alpha(\bar a,\bar b)\) iff
\((\mathcal A,\bar a)\equiv_\alpha(\mathcal B,\bar b)\) by induction on
\(\alpha<\omega_1\), for all tuples. At \(0\) both sides say that the tuples
agree on atomic formulas in their slots. Let \(\alpha=\gamma+1\). If
\(\equiv_{\gamma+1}\) holds, then \(\equiv_\gamma\) holds, so the position is a
partial isomorphism by the case \(0\) and the decrease of the relations in
\cref{thm:p03-morley-approximant-properties}; each move has an answer giving
\(\equiv_\gamma\), hence \(\equiv_\beta\) and, by induction, \(W^\beta\) for
every \(\beta\le\gamma\). By unfolding, \(W^{\gamma+1}\) holds. If
\(W^{\gamma+1}\) holds, unfolding with \(\beta=\gamma\) gives forth and back
into \(W^\gamma={\equiv_\gamma}\), and unfolding of \(W^\gamma\) shows that
\(W^\gamma(\bar a,\bar b)\) itself holds, its conditions being among those of
\(W^{\gamma+1}\); so \(\equiv_{\gamma+1}\) holds. Let \(\lambda\) be a limit.
If \(\equiv_\lambda\) holds, then \(\equiv_0\) gives a partial isomorphism and,
for \(\beta<\lambda\), \(\equiv_{\beta+1}\) gives forth and back into
\(\equiv_\beta=W^\beta\); so \(W^\lambda\) holds. If \(W^\lambda\) holds, then
for each \(\beta<\lambda\) the conditions in the unfolding of \(W^\beta\) are
among those of \(W^\lambda\), so \(W^\beta\), that is \(\equiv_\beta\), holds;
thus \(\equiv_\lambda\) holds. The agreement on infinitary formulas of rank at
most \(\alpha\) is \cref{thm:p03-morley-approximant-properties}.

For finite clocks, let \(V_k(P)\) mean that Duplicator wins \(G_k\) from
\(P\). Winning \(G_{k}\) implies winning \(G_j\) for \(j\le k\): play the same
strategy and stop early; the shorter final position is a restriction of a
longer won one. We show \(V_k=W^k\) by induction. Both are the
partial-isomorphism condition at \(0\). If \(W^{k+1}(P)\), unfolding with
\(\beta=k\) gives forth and back into \(W^k=V_k\), hence \(V_{k+1}(P)\). If
\(V_{k+1}(P)\), then \(P\) is a partial isomorphism and every move has an answer
into \(V_k\subseteq V_\beta=W^\beta\) for all \(\beta\le k\), hence
\(W^{k+1}(P)\).

Finally, if Duplicator wins \(G_\omega\) from \(P\), let \(I\) be a
back-and-forth family containing \(P\)
(\cref{thm:p03-w03-p03-038-omega-comparison}). In \(G^\alpha\), for any
ordinal \(\alpha\), she answers into \(I\); every play is finite and ends in
\(I\), hence in a partial isomorphism.
\end{proof}
\end{annexitem}

\begin{annexitem}{annex:p03-w03-p03-038-long-comparison}{thm:p03-w03-p03-038-long-comparison}{The game of length \(\omega_1\)}
\begin{proof}
\emph{(1).} Let \(\sigma\) win \(G_{\omega_1}\). In \(G_\omega\) Duplicator
uses \(\sigma\) for the first \(\omega\) rounds. Each such play is the
beginning of a play of \(G_{\omega_1}\) following \(\sigma\), in which Spoiler
continues arbitrarily, and all its stages, in particular the finite ones, are
partial isomorphisms. The \(L_{\infty\omega}\) conclusion is
\cref{thm:p03-w03-p03-038-omega-comparison}.

\emph{(2).} Let \(f:\mathcal A\cong\mathcal B\). Duplicator answers \(c\) by
\(f(c)\) and \(d\) by \(f^{-1}(d)\). Every position at every stage is a
restriction of \(f\), hence a partial isomorphism.

\emph{(3).} Since the domains are nonempty and have cardinality at most
\(\aleph_1\), there are surjections \(e_A:\omega_1\to\dom\mathcal A\) and
\(e_B:\omega_1\to\dom\mathcal B\). Each \(\xi<\omega_1\) is uniquely
\(\lambda+i\) with \(\lambda\) zero or a limit and \(i<\omega\). Spoiler
chooses \(e_A(\lambda+i)\) at round \(\lambda+2i\) and \(e_B(\lambda+i)\) at
round \(\lambda+2i+1\); these rounds are below \(\omega_1\). Let Duplicator
follow her winning strategy. The position at stage \(\omega_1\) is a partial
isomorphism; its domain contains every \(e_A(\xi)\) and its range every
\(e_B(\xi)\). Respect for equality makes it the graph of an injective
function, total and onto, that preserves and reflects every relation: an
isomorphism.
\end{proof}
\end{annexitem}

\begin{annexitem}{annex:p03-w03-p03-038-club-facts}{thm:p03-w03-p03-038-club-facts}{Club facts used by the game proofs}
\begin{proof}
\emph{(1).} If \(X\) is a countable set of countable ordinals, \(\sup X=\bigcup X\)
is an ordinal and a countable union of countable sets, hence countable by
countable choice; a countable ordinal is below \(\omega_1\).

\emph{(2).} For \(\alpha<\omega_1\), \(\alpha+\omega\) is a nonzero limit,
countable and above \(\alpha\). If \(\delta\) is a nonzero limit below
\(\omega_1\), then \(\delta\in\mathrm{Lim}\) outright, so closure holds.

\emph{(3).} Let \(C,D\) be clubs. If \(\delta\) is a nonzero limit with
\(\sup(C\cap D\cap\delta)=\delta\), then \(\sup(C\cap\delta)=\delta\) and
\(\sup(D\cap\delta)=\delta\), so \(\delta\in C\cap D\). Given \(\alpha\), let
\(c_0\) be the least member of \(C\) above \(\alpha\), \(d_n\) the least member of
\(D\) above \(c_n\), and \(c_{n+1}\) the least member of \(C\) above \(d_n\). The
supremum \(\delta\) of these countably many ordinals is below \(\omega_1\) by
(1), is a nonzero limit because the sequence increases strictly, and is the
supremum of both \(C\cap\delta\) and \(D\cap\delta\); so \(\delta\in C\cap D\)
and \(\delta>\alpha\). If \(C\) is club and \(D\) is any club, \(C\cap D\) is
club, hence unbounded and nonempty; so \(C\) is stationary. A stationary set
meets every club by \cref{def:p06-club-stationary}.

\emph{(4).} Closure: let \(\delta\) be a nonzero limit that is the supremum of
closure points below it. For \(\gamma<\delta\) choose a closure point
\(\delta'\) with \(\gamma<\delta'<\delta\); then \(g(\gamma)<\delta'<\delta\).
Unboundedness: given \(\alpha\), put \(\gamma_0=\alpha+1\) and
\(\gamma_{n+1}=\sup\{g(\gamma)+1:\gamma\le\gamma_n\}+\gamma_n+1\), which is
countable by (1). The supremum \(\delta\) of the strictly increasing
\(\gamma_n\) is a nonzero limit below \(\omega_1\). If \(\gamma<\delta\) then
\(\gamma\le\gamma_n\) for some \(n\) and
\(g(\gamma)<\gamma_{n+1}<\delta\). So \(\delta\) is a closure point above
\(\alpha\).

\emph{(5).} If \(\delta\) is a nonzero limit with
\(\sup(\mathrm{Lim}'\cap\delta)=\delta\), then \(\sup(\mathrm{Lim}\cap\delta)=\delta\)
since \(\mathrm{Lim}'\subseteq\mathrm{Lim}\), so \(\delta\in\mathrm{Lim}'\).
Given \(\alpha\), let \(\alpha_0=\alpha+\omega\) and
\(\alpha_{n+1}=\alpha_n+\omega\), all in \(\mathrm{Lim}\); their supremum
\(\delta\) is below \(\omega_1\) by (1), above \(\alpha\), and is the supremum
of \(\mathrm{Lim}\cap\delta\), so \(\delta\in\mathrm{Lim}'\). For every
\(\alpha\), the ordinal \(\alpha+\omega\) is in \(\mathrm{Lim}\) but not in
\(\mathrm{Lim}'\): the ordinals strictly between \(\alpha\) and
\(\alpha+\omega\) are successors, so \(\sup(\mathrm{Lim}\cap(\alpha+\omega))\le\alpha\).
Hence \(\mathrm{Lim}\setminus\mathrm{Lim}'\) is unbounded, and it is
nonstationary because it is disjoint from the club \(\mathrm{Lim}'\).
\end{proof}
\end{annexitem}

\begin{annexitem}{annex:p03-w03-p03-038-phi-basic}{thm:p03-w03-p03-038-phi-basic}{Basic properties of \(\Phi(S)\)}
\begin{proof}
\emph{(1).} The lexicographic order of a well-ordered family of linear orders
is linear. Density: let \(x=(\alpha,q)<y=(\beta,q')\). If \(\alpha=\beta\),
then \((\alpha,(q+q')/2)\) lies between them and belongs to the block, since
\(Q_\alpha(S)\) is closed under averages. If \(\alpha<\beta\), then
\((\alpha,q+1)\) lies between them. No greatest element: \((\alpha,q)<(\alpha,q+1)\).
No least element: for \((\alpha,q)\), if \(q-1\in Q_\alpha(S)\) then
\((\alpha,q-1)\) is smaller; otherwise \(\alpha\in S\), so \(\alpha>0\), and
\((0,0)\) is smaller, because \(0\notin S\) gives \(Q_0(S)=\mathbb Q\).
Cardinality: \(\alpha\mapsto(\alpha,0)\) injects \(\omega_1\) into
\(\Phi(S)\), and \(\Phi(S)\subseteq\omega_1\times\mathbb Q\), which injects into
\(\omega_1\times\omega_1\), of cardinality \(\aleph_1\) by
\cref{thm:p01-w02-p01-011-finite-coding}(1). The
Cantor--Schr\"oder--Bernstein theorem gives \(|\Phi(S)|=\aleph_1\).

\emph{(2).} \(\Phi(S)\mathbin{\upharpoonright}\delta\subseteq\delta\times\mathbb Q\),
a product of two countable sets.

\emph{(3).} If \(\delta\in S\), \((\delta,0)\) belongs to the final segment and
lies below every \((\delta,q)\) with \(q\ge0\) and every \((\alpha,q)\) with
\(\alpha>\delta\). If \(\delta\notin S\) and \((\alpha,q)\) is in the final
segment, then \((\delta,q-1)\) is smaller when \(\alpha=\delta\), and
\((\delta,0)\) is smaller when \(\alpha>\delta\); so there is no least element.
\end{proof}
\end{annexitem}

\begin{annexitem}{annex:p03-w03-p03-038-phi-iso}{thm:p03-w03-p03-038-phi-iso}{Isomorphism of \(\Phi(S)\) and \(\Phi(T)\)}
\begin{proof}
\emph{Isomorphism implies nonstationary difference.} Let
\(f:\Phi(S)\cong\Phi(T)\) and
\(C\defeq\{\delta\in\mathrm{Lim}:f[\Phi(S)\mathbin{\upharpoonright}\delta]=
\Phi(T)\mathbin{\upharpoonright}\delta\}\). \(C\) is closed: if \(\delta\) is a
nonzero limit with \(\sup(C\cap\delta)=\delta\), every element of
\(\Phi(S)\mathbin{\upharpoonright}\delta\) or
\(\Phi(T)\mathbin{\upharpoonright}\delta\) has first coordinate below some
\(\gamma\in C\cap\delta\), so both initial segments are unions of the
corresponding segments for \(\gamma\in C\cap\delta\), and \(f\) maps the one
union onto the other. \(C\) is unbounded: given \(\alpha\), let
\(\gamma_0=\alpha+1\), and let \(\gamma_{n+1}>\gamma_n\) exceed every first
coordinate of \(f[\Phi(S)\mathbin{\upharpoonright}\gamma_n]\) and of
\(f^{-1}[\Phi(T)\mathbin{\upharpoonright}\gamma_n]\); these are countable sets
by \cref{thm:p03-w03-p03-038-phi-basic}(2), so such \(\gamma_{n+1}<\omega_1\)
exists by \cref{thm:p03-w03-p03-038-club-facts}(1). For
\(\delta=\sup_n\gamma_n\), \(f[\Phi(S)\mathbin{\upharpoonright}\delta]\subseteq
\Phi(T)\mathbin{\upharpoonright}\delta\) and
\(f^{-1}[\Phi(T)\mathbin{\upharpoonright}\delta]\subseteq
\Phi(S)\mathbin{\upharpoonright}\delta\), so \(\delta\in C\). For \(\delta\in C\),
\(f\) maps the final segment of \(\Phi(S)\) above
\(\Phi(S)\mathbin{\upharpoonright}\delta\) isomorphically onto the corresponding
final segment of \(\Phi(T)\). One has a least element iff the other does, so
\(\delta\in S\) iff \(\delta\in T\) by
\cref{thm:p03-w03-p03-038-phi-basic}(3). Thus the club \(C\) is disjoint
from \(S\mathbin{\triangle}T\), which is therefore nonstationary.

\emph{Nonstationary difference implies isomorphism.} Let \(C\) be a club
disjoint from \(S\mathbin{\triangle}T\), and \(D\defeq C\cap\mathrm{Lim}\), a
club by \cref{thm:p03-w03-p03-038-club-facts}(2),(3). \(D\) is uncountable,
since a countable subset of \(\omega_1\) is bounded by
\cref{thm:p03-w03-p03-038-club-facts}(1); so its increasing enumeration
\(\langle\delta_\xi:\xi<\omega_1\rangle\) has length \(\omega_1\). It is
continuous: for a limit \(\xi\), \(s=\sup_{\eta<\xi}\delta_\eta\) is a nonzero
limit with \(\sup(D\cap s)=s\), hence \(s\in D\); as \(s\) exceeds every
\(\delta_\eta\) with \(\eta<\xi\), \(s\ge\delta_\xi\), and \(s\le\delta_\xi\)
holds trivially. For \(\alpha\ge\delta_0\), the least \(\xi\) with
\(\delta_\xi>\alpha\) is neither \(0\) nor a limit, by continuity, so
\(\delta_\eta\le\alpha<\delta_{\eta+1}\) for a unique \(\eta\).

Hence \(\Phi(S)\) is the ordered sum of the pieces
\(P^S_{-1}\defeq\Phi(S)\mathbin{\upharpoonright}\delta_0\) and
\(P^S_\eta\defeq\{(\alpha,q)\in\Phi(S):\delta_\eta\le\alpha<\delta_{\eta+1}\}\),
\(\eta<\omega_1\), in this order, and likewise for \(T\). Each piece is
countable, nonempty and convex, hence dense. It has no greatest element,
because its upper bound \(\delta_0\) or \(\delta_{\eta+1}\) is a limit, so
\((\alpha,q+1)\) stays in the piece. \(P^S_{-1}\) has no least element, as
\(Q_0(S)=\mathbb Q\) and every element is at least some \((0,q)\). \(P^S_\eta\)
has a least element iff block \(\delta_\eta\) has one, by the argument of
\cref{thm:p03-w03-p03-038-phi-basic}(3), iff \(\delta_\eta\in S\). Since
\(\delta_\eta\in C\), \(\delta_\eta\in S\) iff \(\delta_\eta\in T\).

Corresponding pieces are therefore isomorphic. If neither has a least
element, both are countably infinite dense orders without endpoints, a
nonempty dense order without greatest element being infinite, and
\cref{thm:p03-rn-dlo-uniqueness} makes each isomorphic to \(\mathbb Q\). If
both have least elements, remove them. The remainders are nonempty, dense,
without greatest element, and without least element, since a least element
of the remainder would have nothing between it and the removed point; map
the removed points to each other and the remainders isomorphically. Choose
one isomorphism for each \(\eta\) by the Axiom of Choice. Their union is a
bijection that preserves order inside each piece and between pieces, an
isomorphism \(\Phi(S)\cong\Phi(T)\).
\end{proof}
\end{annexitem}

\begin{annexitem}{annex:p03-w03-p03-038-phi-games}{thm:p03-w03-p03-038-phi-games}{Games on \(\Phi(S)\) and \(\Phi(T)\)}
\begin{proof}
\emph{(1).} Both orders are dense without endpoints
(\cref{thm:p03-w03-p03-038-phi-basic}). For the signature \(\{<\}\) a finite
position is a partial isomorphism iff \(a_i\mapsto b_i\) is an
order-preserving bijection of finite sets. By
\cref{thm:p03-rn-dlo-extension}, such a map extends to any further element
of either domain, so the set of these positions is a back-and-forth family
containing the empty position. Apply
\cref{thm:p03-w03-p03-038-omega-comparison}.

\emph{(2).} By \cref{thm:p03-w03-p03-038-phi-iso} the orders are isomorphic;
apply \cref{thm:p03-w03-p03-038-long-comparison}(2).

\emph{(3).} Fix bijections \(e_S:\omega_1\to\Phi(S)\) and
\(e_T:\omega_1\to\Phi(T)\) (\cref{thm:p03-w03-p03-038-phi-basic}(1)). Write
\(P_\xi\) for the position at stage \(\xi\). Spoiler's strategy at round
\(\xi\):
\begin{enumerate}
\item[(a)] if \(\xi\in\mathrm{Lim}\cap(S\mathbin{\triangle}T)\), and \(P_\xi\)
      is a partial isomorphism with domain \(\Phi(S)\mathbin{\upharpoonright}\xi\)
      and range \(\Phi(T)\mathbin{\upharpoonright}\xi\), \emph{attack}: choose
      \((\xi,0)\) in \(\Phi(S)\) if \(\xi\in S\), in \(\Phi(T)\) if \(\xi\in T\);
\item[(b)] if \(\xi=\zeta+1\) and round \(\zeta\) was an attack, answered by
      \(v\): if \(v\) lies in the final segment of the other order above stage
      \(\zeta\), say \(v=(\beta,y)\) with \(\beta\ge\zeta\), choose \(v'=(\zeta,y-1)\)
      if \(\beta=\zeta\) and \(v'=(\zeta,0)\) if \(\beta>\zeta\), in that order;
      otherwise choose anything;
\item[(c)] otherwise write \(\xi=\lambda+n\) with \(\lambda\) zero or a limit,
      and choose \(e_S(\lambda+i)\) in \(\Phi(S)\) if \(n=2i\), or
      \(e_T(\lambda+i)\) in \(\Phi(T)\) if \(n=2i+1\).
\end{enumerate}

\emph{Claim A: an attack wins.} Suppose round \(\zeta\) is an attack with
\(\zeta\in S\setminus T\); the case \(\zeta\in T\setminus S\) exchanges the
orders. Spoiler chose \(u=(\zeta,0)\), which is not in the domain
\(\Phi(S)\mathbin{\upharpoonright}\zeta\) of \(P_\zeta\). If the answer \(v\)
lies in the range \(\Phi(T)\mathbin{\upharpoonright}\zeta\), it is already
paired with an element different from \(u\), and \(P_{\zeta+1}\) fails to
respect equality. Otherwise \(v\) is in the final segment of \(\Phi(T)\) above
\(\zeta\), which has no least element because \(\zeta\notin T\); the point
\(v'\) of (b) lies in that final segment, below \(v\). Let \(u'\) answer
\(v'\). If \(P_{\zeta+2}\) were a partial isomorphism, then \(u'<u\) because
\(v'<v\), and \(u'\) would exceed every element of
\(\Phi(S)\mathbin{\upharpoonright}\zeta\) because \(v'\) exceeds every element of
\(\Phi(T)\mathbin{\upharpoonright}\zeta\). But \(u\) is the least element above
\(\Phi(S)\mathbin{\upharpoonright}\zeta\) by
\cref{thm:p03-w03-p03-038-phi-basic}(3). So Duplicator loses at stage
\(\zeta+1\) or \(\zeta+2\).

\emph{Claim B: an attack occurs unless Duplicator has already lost.} Suppose
a play follows Spoiler's strategy and every position in it is a partial
isomorphism. By Claim A no attack occurs, so every round is of type (c). Let
Spoiler's move at round \(\eta\) be \(w_\eta\) and the pair added at round
\(\eta\) be \((u_\eta,v_\eta)\). For \(\gamma<\omega_1\) let \(g(\gamma)\) be the
least ordinal above \(\gamma\) that exceeds every round at which an element of
\(\Phi(S)\mathbin{\upharpoonright}\gamma\) or \(\Phi(T)\mathbin{\upharpoonright}\gamma\)
is chosen by Spoiler and every first coordinate of \(u_\eta\) and \(v_\eta\)
for \(\eta<\gamma\). These are countably many countable ordinals, because the
segments are countable and each element is chosen at the single round fixed
by (c); so \(g(\gamma)<\omega_1\) by \cref{thm:p03-w03-p03-038-club-facts}(1).
Let \(\delta\) be a closure point of \(g\). Each \(x=(\alpha,q)\) in
\(\Phi(S)\mathbin{\upharpoonright}\delta\) lies in
\(\Phi(S)\mathbin{\upharpoonright}(\alpha+1)\) with \(\alpha+1<\delta\), so it is
chosen by Spoiler before round \(g(\alpha+1)<\delta\) and lies in the domain of
\(P_\delta\); similarly \(\Phi(T)\mathbin{\upharpoonright}\delta\) lies in its
range. Each \(u_\eta\) with \(\eta<\delta\) has first coordinate below
\(g(\eta+1)<\delta\), and likewise each \(v_\eta\). Hence \(P_\delta\) is a
partial isomorphism with domain \(\Phi(S)\mathbin{\upharpoonright}\delta\) and
range \(\Phi(T)\mathbin{\upharpoonright}\delta\). The closure points form a club
contained in \(\mathrm{Lim}\) (\cref{thm:p03-w03-p03-038-club-facts}(4)), and
\(S\mathbin{\triangle}T\) is stationary, so some closure point \(\delta\) lies in
\(S\mathbin{\triangle}T\). At round \(\delta\) clause (a) applies and Spoiler
attacks, contradicting the absence of attacks.

By Claims A and B every play following the strategy contains a stage whose
position is not a partial isomorphism. Spoiler wins.
\end{proof}
\end{annexitem}

\begin{annexitem}{annex:p03-w03-p03-038-phi-inexpressible}{thm:p03-w03-p03-038-phi-inexpressible}{An inexpressibility result from a stationary set}
\begin{proof}
\(\varnothing\mathbin{\triangle}\mathrm{Lim}=\mathrm{Lim}\) is club by
\cref{thm:p03-w03-p03-038-club-facts}(2), hence stationary by (3). By
\cref{thm:p03-w03-p03-038-phi-games}(1) the two orders satisfy the same
\(L_{\infty\omega}\)-sentences; by
\cref{thm:p03-w03-p03-038-phi-iso} they are not isomorphic; by
\cref{thm:p03-w03-p03-038-phi-games}(3) Spoiler wins the game of length
\(\omega_1\). If an \(L_{\infty\omega}\)-sentence \(\sigma\) held exactly in the
orders isomorphic to \(\Phi(\mathrm{Lim})\), then
\(\Phi(\mathrm{Lim})\models\sigma\), hence \(\Phi(\varnothing)\models\sigma\),
hence \(\Phi(\varnothing)\cong\Phi(\mathrm{Lim})\), a contradiction. Both
orders are dense without endpoints and of cardinality \(\aleph_1\)
(\cref{thm:p03-w03-p03-038-phi-basic}), so the restriction to that class
does not help.
\end{proof}
\end{annexitem}
